\documentclass[11pt,a4paper,oneside]{scrreprt}

\usepackage[utf8]{inputenc}
\usepackage[T1]{fontenc}
\usepackage[ngerman]{babel}
\usepackage{lmodern}
\usepackage{microtype}
\usepackage{geometry}
\usepackage{booktabs}
\usepackage{longtable}
\usepackage{tabularx}
\usepackage{array}
\usepackage{enumitem}
\usepackage{csquotes}
\usepackage{xcolor}
\usepackage{graphicx}
\usepackage{tikz}
\usepackage{rotating}
\usepackage{hyperref}
\usepackage[backend=biber,style=ieee]{biblatex}

\usetikzlibrary{arrows.meta,positioning,fit,calc,shapes.geometric,shapes.misc,backgrounds,matrix}

\geometry{left=28mm,right=28mm,top=28mm,bottom=30mm}
\addbibresource{bibliography/references.bib}

\hypersetup{
  colorlinks=true,
  linkcolor=black,
  citecolor=black,
  urlcolor=blue,
  pdftitle={MBSE-gestuetzte Entwicklung eines mobilen Materialtransportroboters},
  pdfauthor={Projektteam}
}

\newcommand{\projectversion}{0.3}
\newcommand{\docstatus}{Arbeitsstand}
\newcommand{\sysml}{SysML v2}
\newcommand{\ros}{ROS2}

\newcolumntype{L}[1]{>{\raggedright\arraybackslash}p{#1}}
\newcolumntype{Y}{>{\raggedright\arraybackslash}X}

\setlist[itemize]{topsep=3pt,itemsep=2pt,parsep=0pt}
\setlist[enumerate]{topsep=3pt,itemsep=2pt,parsep=0pt}
\renewcommand{\arraystretch}{1.18}

\newenvironment{decisionbox}[1]
  {\par\medskip\noindent\textbf{#1}\par\smallskip\noindent}
  {\par\medskip}

% Modellgebundene Tabellenzeilen aus tables/generated/ (erzeugt mit tools/generate_docs.py).
% Das primitive \@@input funktioniert auch innerhalb von longtable.
\makeatletter
\newcommand{\modeltable}[1]{\@@input tables/generated/#1.tex }
\newcommand{\modeltext}[1]{\@@input tables/generated/#1.tex }
\makeatother


\title{MBSE-gestuetzte Entwicklung eines mobilen Materialtransportroboters}
\subtitle{Begleitende LaTeX-Dokumentation fuer SysML v2, MVP, Architektur, ROS2-Mapping, Verifikation und Traceability}
\author{Projektteam}
\date{Stand: \today \\ Version \projectversion}

\begin{document}
\pagenumbering{roman}
\maketitle

\begin{abstract}
Dieses Dokument begleitet die schrittweise Entwicklung eines mobilen Indoor-Roboters zum Transport kleiner Materialkisten. Es dient als lebendes Engineering-Dokument und als Grundlage fuer eine spaetere publikationsfaehige Overleaf-Fassung. Der Entwicklungsprozess folgt einem MBSE-Ansatz: Stakeholder-Beduerfnisse werden in Anforderungen ueberfuehrt, Anforderungen werden mit Use Cases, Funktionen, Architektur, Schnittstellen, ROS2-Elementen und Verifikationsfaellen verbunden. Der MVP wird dabei nicht als Wegwerfprototyp betrachtet, sondern als erster durchgaengiger Ausschnitt des Zielsystems.
\end{abstract}

\tableofcontents
\listoffigures
\listoftables
\clearpage
\pagenumbering{arabic}

\chapter{Einleitung}

\section{Projektanlass}

Ziel des Projekts ist die Entwicklung eines autonomen mobilen Indoor-Roboters, der kleine Materialkisten zwischen festgelegten Standorten transportiert. Die Standorte sind in einer internen Karte hinterlegt. Eine grobe Karte des Bewegungsraums liegt zu Projektbeginn vor, wird jedoch waehrend des Betriebs aktualisiert, damit temporaere und dauerhafte Hindernisse in der Navigation beruecksichtigt werden koennen.

Das Projekt verfolgt zwei miteinander gekoppelte Ziele:
\begin{enumerate}
  \item Entwicklung eines funktionsfaehigen mobilen Materialtransportroboters.
  \item Aufbau einer nachvollziehbaren MBSE-Entwicklungskette, in der Anforderungen, \sysml-Modell, Architektur, Hardware, Software, \ros-Mapping, Tests und Dokumentation konsistent miteinander verbunden bleiben.
\end{enumerate}

\section{Rolle dieses Dokuments}

Dieses Dokument wird parallel zur technischen Entwicklung gepflegt. Es ist weder ein nachtraegliches Protokoll noch ein reines Pflichtenheft. Es ist die zentrale textuelle Projektdokumentation fuer:
\begin{itemize}
  \item Entwicklungsschritte und Iterationen,
  \item Annahmen, Randbedingungen und offene Punkte,
  \item Stakeholder, Beduerfnisse, Requirements und Use Cases,
  \item MVP-Abgrenzung gegenueber dem Zielsystem,
  \item Architekturentscheidungen und betrachtete Alternativen,
  \item Hardware-/Software-Allokation,
  \item Schnittstellen und \ros-Mapping,
  \item Verifikation, Validierung und Traceability.
\end{itemize}

Die Struktur ist modular aufgebaut, damit einzelne Kapitel in Overleaf gezielt erweitert, reviewed und spaeter in eine publikationsfaehige Fassung ueberfuehrt werden koennen.

\section{Quellen und Ausgangsdaten}

Die initiale Fassung basiert auf:
\begin{itemize}
  \item dem bisherigen Projektchat zur Rolle von KI als Entwicklungspartner,
  \item der Projektbeschreibung zum mobilen Materialtransportroboter,
  \item dem bereitgestellten Lehrbuch \emph{SysML v2, MBSE und ROS2 fuer Einsteiger} \cite{becker2026sysml},
  \item den daraus abgeleiteten MBSE-Arbeitsregeln fuer Anforderungen, Modellierung, Architektur, Implementierung und Nachweis.
\end{itemize}

Das Lehrbuch wird als methodischer Rahmen verwendet. Es betont insbesondere, dass technische Loesungen nicht vor dem Problemverstaendnis festgelegt werden sollen, dass Stakeholder, Systemgrenze und Kontext frueh zu klaeren sind und dass Tests und Nachweise nicht erst am Projektende geplant werden duerfen \cite[Kap.~1]{becker2026sysml}. Fuer das vorliegende Projekt wird diese Methodik um einen expliziten MVP-Schnitt innerhalb des Gesamtmodells erweitert.

\section{Dokumentationsregel}

Jeder wesentliche Entwicklungsschritt wird nach demselben Muster dokumentiert:
\begin{itemize}
  \item Ausgangslage und Problem,
  \item Annahmen,
  \item betroffene Stakeholder,
  \item abgeleitete Anforderungen,
  \item betrachtete Alternativen,
  \item Entscheidung,
  \item Begruendung,
  \item Auswirkungen auf Zielsystem und MVP,
  \item offene Punkte,
  \item zugehoerige \sysml-Elemente,
  \item Verifikations- oder Validierungsidee,
  \item Aenderungshistorie.
\end{itemize}

Diese Regel verhindert, dass Architekturwissen nur implizit im Chat, in Code, in Zeichnungen oder in Einzelentscheidungen verbleibt.

\chapter{Methodik und Entwicklungsablauf}

\section{Grundprinzip}

Die Entwicklung folgt einem inkrementellen MBSE-Prozess. Zuerst werden Zweck, Kontext und Systemgrenze verstanden. Danach werden Stakeholder-Beduerfnisse in testbare Anforderungen ueberfuehrt. Aus Anforderungen entstehen Use Cases, Funktionen, logische und physische Architekturen, Schnittstellen, \ros-Elemente und Verifikationsfaelle.

Das Lehrbuch beschreibt MBSE nicht als Produktion moeglichst vieler Diagramme, sondern als systematische Verknuepfung von Informationen: Stakeholder-Bedarf, Anforderung, Funktion, Systemelement und Testfall muessen nachvollziehbar miteinander verbunden sein \cite[Kap.~2]{becker2026sysml}. Diese Traceability ist fuer das Roboterprojekt besonders wichtig, weil Aenderungen an Navigation, Safety, Sensorik oder Energieversorgung mehrere Entwicklungsbereiche gleichzeitig betreffen koennen.

\section{Iterativer Ablauf}

Der Entwicklungsablauf wird in Iterationen organisiert. Jede Iteration erzeugt konkrete Artefakte und aktualisiert dieses Dokument.

\begin{figure}[htbp]
  \centering
  \resizebox{\textwidth}{!}{\begin{tikzpicture}[
  node distance=9mm,
  step/.style={draw, rounded corners=2pt, align=center, minimum width=31mm, minimum height=8mm, fill=gray!8},
  artifact/.style={draw, rounded corners=2pt, align=center, minimum width=31mm, minimum height=8mm, fill=blue!6},
  arrow/.style={-{Latex[length=2mm]}, thick},
  feedback/.style={-{Latex[length=2mm]}, thick, dashed}
]
\node[step] (goal) {Systemzweck\\und Kontext};
\node[step, right=of goal] (needs) {Stakeholder\\Needs};
\node[step, right=of needs] (req) {Requirements\\und Use Cases};
\node[step, right=of req] (func) {Funktionen\\und Verhalten};
\node[step, below=of func] (arch) {Architektur\\und Schnittstellen};
\node[step, left=of arch] (mvp) {MVP-Schnitt\\im Zielsystem};
\node[step, left=of mvp] (impl) {HW/SW, ROS2\\und Mechanik};
\node[artifact, left=of impl] (verify) {Verifikation\\und Traceability};

\draw[arrow] (goal) -- (needs);
\draw[arrow] (needs) -- (req);
\draw[arrow] (req) -- (func);
\draw[arrow] (func) -- (arch);
\draw[arrow] (arch) -- (mvp);
\draw[arrow] (mvp) -- (impl);
\draw[arrow] (impl) -- (verify);
\draw[feedback] (verify.north) to[out=90,in=270] node[above, align=center, font=\small] {Erkenntnisse aktualisieren\\Modell und Entscheidungen} (needs.south);
\end{tikzpicture}
}
  \caption{Inkrementeller Entwicklungsablauf vom Systemzweck bis zur Verifikation}
  \label{fig:development-flow}
\end{figure}

\begin{longtable}{L{0.16\textwidth} L{0.30\textwidth} L{0.42\textwidth}}
\toprule
Iteration & Schwerpunkt & Ergebnisartefakte \\
\midrule
A & Systemdefinition & Zielsystem, Systemgrenze, Einsatzumgebung, Stakeholder, Annahmen, erste Risiken \\
B & Requirements und Use Cases & Stakeholder Needs, Requirements, Use Cases, Akzeptanzkriterien \\
C & MVP-Definition & MVP-Scope, Nicht-MVP, Erfolgskriterien, MVP-Traceability \\
D & Systemarchitektur & Funktionen, logische Architektur, physische Architektur, Schnittstellen, Betriebsmodi \\
E & Implementierungsarchitektur & HW/SW-Allokation, Rechnerarchitektur, Sensorik, Elektronik, \ros-Mapping \\
F & Verifikation & Testfaelle, Simulation, Integrationsstrategie, Systemtests, Traceability-Matrix \\
\bottomrule
\end{longtable}

\section{Single Source of Truth}

Das Gesamtmodell bleibt die fachliche Referenz. Der MVP wird nicht als separates Modell erstellt, sondern als markierter Ausschnitt des Zielsystems. Jedes relevante Element erhaelt daher mindestens:
\begin{itemize}
  \item eine stabile ID,
  \item einen Status: \emph{MVP}, \emph{Release 1}, \emph{Future} oder \emph{Offen},
  \item eine Prioritaet: \emph{Must}, \emph{Should}, \emph{Could},
  \item eine Quelle oder Begruendung,
  \item eine geplante Nachweismethode.
\end{itemize}

\section{Rolle von KI im Prozess}

KI wird als Engineering-Partner eingesetzt, nicht als alleinige Freigabeinstanz. Sie unterstuetzt beim Strukturieren, Generieren und Pruefen von Artefakten. Formale Modellpruefung, Berechnungen, Simulation, Compiler, Tests, CAD-/EDA-Pruefungen und menschliche Reviews bleiben entscheidend.

Die Projektregel lautet:
\begin{quote}
KI darf Entwicklungsartefakte vorschlagen und konsistent halten. Technische Gueltigkeit entsteht erst durch nachvollziehbare Anforderungen, begruendete Entscheidungen, formale Modellierung, Berechnung, Simulation, Test und Review.
\end{quote}

\section{Qualitaetskriterien fuer Artefakte}

Ein Artefakt gilt erst dann als belastbar, wenn folgende Fragen beantwortet sind:
\begin{itemize}
  \item Welche Anforderung oder welches Stakeholder-Beduerfnis begruendet das Artefakt?
  \item Welche Alternative wurde betrachtet?
  \item Warum wurde die aktuelle Loesung gewaehlt?
  \item Welche Schnittstellen, Tests und Risiken entstehen daraus?
  \item Wie wirkt sich eine spaetere Aenderung auf andere Artefakte aus?
\end{itemize}

\section{Entwicklungsprinzipien aus Iteration A}

Aus der Vervollstaendigung von Iteration A werden folgende Entwicklungsprinzipien als Projektregeln uebernommen:

\begin{longtable}{L{0.14\textwidth} L{0.28\textwidth} L{0.44\textwidth}}
\toprule
ID & Prinzip & Bedeutung \\
\midrule
DP-01 & Requirement vor Loesung & Technologien werden erst festgelegt, wenn Anforderungen sie begruenden. \\
DP-02 & Gesamtmodell vor MVP & Das Gesamtprojekt wird modelliert; der MVP ist eine Teilmenge davon. \\
DP-03 & Keine Wegwerfarchitektur & MVP-Komponenten sollen moeglichst in die spaetere Architektur ueberfuehrbar sein. \\
DP-04 & Traceability & Needs, Requirements, Use Cases, Funktionen, Architektur, Implementierung und Verifikation werden verknuepft. \\
DP-05 & Entscheidungen dokumentieren & Groessere Architekturentscheidungen erhalten einen Decision Record. \\
DP-06 & Unsicherheit sichtbar machen & Requirements, Decisions, Assumptions, Open Issues und Risks werden getrennt dokumentiert. \\
DP-07 & Schrittweise Verifikation & Nicht erst der fertige Roboter wird getestet; Teilmodelle, Simulationen, Baugruppen und Software werden frueh geprueft. \\
\bottomrule
\end{longtable}

\section{Artefakttypen und Umgang mit Unsicherheit}

Fuer die weitere Entwicklung wird strikt unterschieden zwischen:
\begin{itemize}
  \item \emph{Requirement}: nachzuweisende Forderung an das System,
  \item \emph{Decision}: getroffene und begruendete Architektur- oder Prozessentscheidung,
  \item \emph{Assumption}: vorlaeufige Annahme, die spaeter bestaetigt oder ersetzt werden muss,
  \item \emph{Open Issue}: offene Frage mit Auswirkung auf Anforderungen, Architektur oder Tests,
  \item \emph{Risk}: moegliches Ereignis oder technische Unsicherheit mit negativer Auswirkung.
\end{itemize}

Diese Trennung ist wichtig, damit offene Fragen nicht versehentlich als Anforderungen erscheinen und Loesungsideen nicht zu frueh als feste Randbedingungen behandelt werden.

\chapter{Systemdefinition Baseline A.1}
\label{chap:systemdefinition}

\section{Zweck dieses Kapitels}

Dieses Kapitel bildet die fachliche Baseline der Iteration A. Es beschreibt Zielsystem, Systemkontext, Systemgrenze, Betriebsbedingungen, Entscheidungen, Annahmen, Nicht-Ziele, Schnittstellen und offene Punkte. Nach der Methodik des Lehrbuchs wird zuerst vom Problem, den Stakeholdern, der Umgebung und den nachzuweisenden Wirkungen ausgegangen. Konkrete Loesungen wie Sensorfabrikate, Fahrwerkskonzept, Rechnerplattform oder konkrete \ros-Pakete werden erst spaeter entschieden \cite[Kap.~1--2]{becker2026sysml}.

\section{Arbeitstitel und Kurzname}

Der Arbeitstitel des Systems lautet:
\begin{quote}
Autonomer mobiler Materialtransportroboter.
\end{quote}

Als Kurzname wird fuer die weitere Dokumentation \emph{Mobile Material Robot} verwendet, abgekuerzt \emph{MMR}. Der Name dient der eindeutigen Referenz im Projekt und stellt noch keine Produkt- oder Markenentscheidung dar.

\section{Problemstellung}

In einem produktionsaehnlichen Bereich einer Schule sollen Werkstuecktraeger eines Transfersystems zwischen definierten Standorten transportiert werden. Der Fahrbereich umfasst mehrere Raeume und einen verbindenden Flur. Der Roboter muss offene Tueren durchfahren koennen und in einer Umgebung arbeiten, die nicht vollstaendig statisch ist. Personen, abgestellte Gegenstaende oder andere Hindernisse koennen sich im Fahrbereich befinden oder ihre Position veraendern.

Eine Grundkarte des Betriebsbereichs ist vorhanden. Diese Karte reicht fuer den Betrieb allein nicht aus, weil der Roboter aktuelle Wahrnehmungsdaten beruecksichtigen muss. Das System muss daher bekannte Umgebungsinformation mit aktueller Wahrnehmung verbinden und daraus ein nutzbares Navigationsmodell ableiten.

\section{Systemzweck}

Der Systemzweck wird fuer Baseline A.1 wie folgt festgehalten:
\begin{quote}
Der Mobile Material Robot transportiert Werkstuecktraeger eines Transfersystems autonom zwischen definierten Standorten innerhalb eines produktionsaehnlichen Indoor-Demonstrationsbereichs einer Schule.
\end{quote}

Das System soll dazu:
\begin{itemize}
  \item Transportauftraege entgegennehmen,
  \item Start- und Zielstandorte kennen,
  \item die eigene Position bestimmen,
  \item die Umgebung wahrnehmen,
  \item Aenderungen gegenueber der bekannten Umgebung erkennen,
  \item geeignete Fahrwege bestimmen,
  \item Hindernisse beruecksichtigen,
  \item autonom zum Ziel fahren,
  \item den Transportzustand melden,
  \item auf Fehler oder nicht loesbare Situationen kontrolliert reagieren.
\end{itemize}

Der didaktische Zweck ist ausdruecklich Teil des Projekts. Der Roboter dient nicht nur als Transportmittel, sondern auch als Demonstrator fuer Systems Engineering, \sysml, mobile Robotik, Sensorik, Navigation, Elektronik, Softwareentwicklung, \ros, Verifikation und Test.

\section{System of Interest}

Das \emph{System of Interest} ist der mobile Materialtransportroboter einschliesslich der fuer seinen autonomen Betrieb notwendigen Onboard-Hardware und Onboard-Software. Dazu gehoeren auf Systemebene insbesondere:
\begin{itemize}
  \item mechanische Roboterplattform,
  \item Fahrwerk und Antrieb,
  \item Energieversorgung auf dem Roboter,
  \item Onboard-Sensorik,
  \item Onboard-Rechentechnik und eingebettete Controller,
  \item lokale Software fuer Mission, Navigation, Wahrnehmung, Diagnose und Status,
  \item Kommunikationsschnittstellen,
  \item Transportaufnahme fuer den Werkstuecktraeger,
  \item Bedien- und Statusfunktionen,
  \item Sicherheitsfunktionen.
\end{itemize}

Diese Aufteilung ist eine fachliche Systemstruktur, noch keine endgueltige Architektur. Sie dient dazu, spaetere Requirements, Schnittstellen und Tests sauber zu verorten.

\section{Systemkontext und Systemgrenze}

Abbildung~\ref{fig:system-context-boundary} zeigt den Systemkontext auf Baseline-Niveau. Die Systemgrenze verlaeuft um den Roboter selbst. Externe Akteure und Systeme bleiben sichtbar, weil sie Anforderungen und Schnittstellen erzeugen.

\begin{figure}[htbp]
  \centering
  \resizebox{\textwidth}{!}{\begin{tikzpicture}[
  actor/.style={draw, rounded corners=2pt, align=center, minimum width=29mm, minimum height=8mm, fill=gray!8},
  system/.style={draw, rounded corners=3pt, align=center, minimum width=58mm, minimum height=44mm, fill=blue!5, thick},
  inside/.style={align=center, font=\small},
  arrow/.style={-{Latex[length=2mm]}, thick}
]
\node[system] (robot) {Mobile Material Robot\\[1mm]
  \begin{tabular}{c}
  Mission / Auftrag\\
  Navigation / Lokalisierung\\
  Perception / Karte\\
  Motion / Energie\\
  Transport Interface\\
  Safety / Diagnose
  \end{tabular}};

\node[actor, above=17mm of robot] (web) {Web-App\\Browser};
\node[actor, left=25mm of robot] (carrier) {Werkstuecktraeger\\Payload};
\node[actor, below left=16mm and 18mm of robot] (transfer) {externes\\Transfersystem};
\node[actor, right=24mm of robot] (env) {Fahrumgebung\\Raeume, Flur, Tueren};
\node[actor, below right=16mm and 18mm of robot] (people) {Personen und\\Hindernisse};
\node[actor, below=18mm of robot] (charge) {Ladestation};

\draw[arrow] (web) -- node[right, font=\small] {Transportauftrag\\Status} (robot);
\draw[arrow] (carrier) -- node[above, font=\small] {mechanische\\Auflage} (robot);
\draw[arrow] (transfer) -- node[below left, font=\small] {Uebergabe\\spaeter} (robot);
\draw[arrow] (robot) -- node[above, font=\small] {Bewegung\\Wahrnehmung} (env);
\draw[arrow] (people) -- node[below right, font=\small] {dynamische\\Ereignisse} (robot);
\draw[arrow] (robot) -- node[right, font=\small] {Rueckkehr\\Energie} (charge);
\end{tikzpicture}
}
  \caption{Systemkontext und vorlaeufige Systemgrenze des Mobile Material Robot}
  \label{fig:system-context-boundary}
\end{figure}

\subsection{Innerhalb der Systemgrenze}

\begin{longtable}{L{0.30\textwidth} L{0.58\textwidth}}
\toprule
Element & Begruendung \\
\midrule
Mobile Roboterplattform & Physischer Traeger aller Transport-, Navigations- und Energiekomponenten. \\
Fahrwerk, Antrieb und Motion Control & Notwendig zur Ausfuehrung geplanter Bewegungen. \\
Energieversorgung auf dem Roboter & Erforderlich fuer autonome Missionen und fuer sichere Reaktionen bei niedrigem Ladezustand. \\
Onboard-Sensorik & Grundlage fuer Lokalisierung, Kartierung, Hinderniserkennung und Safety-Funktionen. \\
Onboard-Rechner und Controller & Ausfuehrung von Mission, Navigation, Wahrnehmung, Kommunikation, Diagnose und hardwarenaher Regelung. \\
Transportaufnahme & Mechanische Schnittstelle zum Werkstuecktraeger. \\
Onboard-Software & Realisiert Missionslogik, lokale Bedienung, Status, Diagnose, Navigation, Perception und Safety-Reaktionen. \\
Kommunikationsschnittstellen & Notwendig fuer Web-App, Statusmeldungen und spaetere Integration. \\
\bottomrule
\end{longtable}

\subsection{Ausserhalb der Systemgrenze}

\begin{longtable}{L{0.30\textwidth} L{0.58\textwidth}}
\toprule
Externer Bestandteil & Schnittstellenrelevanz \\
\midrule
Bediener und Betreiber & Erteilen Auftraege, loesen Betriebsbefehle aus, erhalten Status und Fehlerdiagnosen. \\
Web-App und Browser & Erste Auftragsquelle und Bedienoberflaeche; technische Ausfuehrung wird spaeter festgelegt. \\
Werkstuecktraeger & Wird transportiert, ist aber selbst nicht Teil des Roboters. \\
Externes Transfersystem & Kann den Werkstuecktraeger perspektivisch auf den Roboter uebergeben. \\
Raeume, Flur, offene Tueren & Betriebsumgebung und Randbedingung fuer Navigation, Dimensionierung und Test. \\
Personen im Fahrbereich & Externe Akteure mit Sicherheitsrelevanz. \\
Bewegliche Gegenstaende und Hindernisse & Externe Objekte, die erkannt, lokal beruecksichtigt oder gemeldet werden muessen. \\
Ladestation & Ziel fuer Rueckkehrfunktion; automatisches Docking und Laden bleiben fuer den MVP noch offen. \\
IT- und Netzwerkinfrastruktur & Moegliche Grundlage fuer Web-App, Updates, Logging und spaetere Integration. \\
\bottomrule
\end{longtable}

\section{Transportgut und mechanische Schnittstelle}

Das Transportgut ist ein Werkstuecktraeger mit Standardmassen eines Transfersystems. Er wird in der ersten Version nicht aktiv vom Roboter aufgenommen. Vorgesehen sind zwei Beladungsarten:
\begin{itemize}
  \item ein Mensch setzt den Werkstuecktraeger auf den Roboter,
  \item ein externes Transfersystem uebergibt den Werkstuecktraeger an den Roboter.
\end{itemize}

Damit liegt die aktive Be- und Entladung ausserhalb des MVP-Kerns. Der Roboter benoetigt jedoch eine definierte mechanische Transportaufnahme. In Iteration B muessen Breite, Laenge, Hoehe, Eigenmasse, maximale Werkstueckmasse, Gesamtmasse der Payload und Schwerpunktbereich festgelegt werden.

\section{Einsatzumgebung}

\begin{longtable}{L{0.28\textwidth} L{0.56\textwidth}}
\toprule
Merkmal & Festlegung Baseline A.1 \\
\midrule
Einsatzgebiet & Produktionsaehnlicher Schulbereich. \\
Betriebsart & Demonstrator und Lehrsystem. \\
Umgebung & Indoor. \\
Raeumlicher Umfang & Mehrere Raeume und ein verbindender Flur. \\
Tueren & Offene Tueren muessen durchfahren werden. \\
Boden & Weitgehend ebener, fester Boden; kleine Schwellen bleiben zu pruefen. \\
Personen & Koennen sich im Fahrbereich befinden. \\
Gegenstaende & Koennen Fahrwege temporaer blockieren. \\
Grundkarte & Vorhanden. \\
Aktuelle Umgebung & Wird waehrend des Betriebs erfasst. \\
Transportziele & Definierte Standorte. \\
\bottomrule
\end{longtable}

\section{Karten- und Navigationslogik}

Die Kartenlogik wird bewusst in mehrere Informationsarten getrennt. Der Roboter soll nicht blind einer statischen Karte folgen und auch nicht unkontrolliert jede Beobachtung dauerhaft in die Langzeitkarte schreiben.

\begin{figure}[htbp]
  \centering
  \resizebox{\textwidth}{!}{\begin{tikzpicture}[
  block/.style={draw, rounded corners=2pt, align=center, minimum width=38mm, minimum height=9mm, fill=gray!8},
  data/.style={draw, rounded corners=2pt, align=center, minimum width=38mm, minimum height=9mm, fill=blue!6},
  decision/.style={draw, diamond, aspect=2.2, align=center, inner sep=1pt, fill=yellow!12},
  arrow/.style={-{Latex[length=2mm]}, thick}
]
\node[data] (base) {Basisinformationen\\Langzeitkarte};
\node[data, below=9mm of base] (sensor) {aktuelle\\Sensorinformationen};
\coordinate (inputmid) at ($(base)!0.5!(sensor)$);
\node[block, right=24mm of inputmid] (fusion) {Umgebungsmodell\\zur Laufzeit};
\node[decision, right=22mm of fusion] (kind) {Aenderung\\temporaer?};
\node[block, above right=8mm and 18mm of kind] (local) {lokale Hindernisse\\beruecksichtigen};
\node[block, below right=8mm and 18mm of kind] (review) {Kartenupdate\\markieren / freigeben};
\node[block, right=25mm of local] (nav) {Navigation\\und Pfadplanung};
\node[block, below=13mm of nav] (motion) {Fahrentscheidung\\und Bewegung};

\draw[arrow] (base) -- (fusion);
\draw[arrow] (sensor) -- (fusion);
\draw[arrow] (fusion) -- (kind);
\draw[arrow] (kind) -- node[above, font=\small] {ja} (local);
\draw[arrow] (kind) -- node[below, font=\small] {nein / unklar} (review);
\draw[arrow] (local) -- (nav);
\draw[arrow] (review) -| (nav);
\draw[arrow] (nav) -- (motion);
\end{tikzpicture}
}
  \caption{Trennung von Grundkarte, aktueller Wahrnehmung und Navigationsentscheidung}
  \label{fig:navigation-logic}
\end{figure}

\begin{longtable}{L{0.24\textwidth} L{0.42\textwidth} L{0.22\textwidth}}
\toprule
Informationsart & Inhalt & Offene Entscheidung \\
\midrule
Basisinformationen & Waende, Raeume, Flure, feste Durchgaenge, bekannte Transportstationen und dauerhaft gesperrte Bereiche. & Erstellung, Freigabe und Versionierung der Grundkarte. \\
Aktuelle Umgebungsinformationen & Personen, abgestellte Kisten, Fahrzeuge, Moebel, temporaere Sperren und andere Hindernisse. & Gueltigkeitsdauer und Umgang mit Sensorunsicherheit. \\
Dauerhaft erkannte Aenderungen & Veraenderungen, die fuer die Langzeitkarte relevant sein koennten. & Menschliche Freigabe oder automatische Uebernahme. \\
\bottomrule
\end{longtable}

\section{Leistungs- und Betriebsrandbedingungen}

\begin{longtable}{L{0.22\textwidth} L{0.38\textwidth} L{0.26\textwidth}}
\toprule
Thema & Festlegung Baseline A.1 & Bemerkung \\
\midrule
Nominalgeschwindigkeit MVP & ca. $0{,}5\,\mathrm{m/s}$, entsprechend ca. $1{,}8\,\mathrm{km/h}$ & Bewusst niedrig fuer Demonstration, Navigation, Tests und Unterricht. \\
Spaetere Geschwindigkeit & Hoehere Geschwindigkeit soll perspektivisch moeglich bleiben. & Future Requirement; keine sofortige Auslegung auf Maximalgeschwindigkeit. \\
Betriebsdauer erste Version & ca. 2 h & Deckt einen Unterrichtsblock von 90 min mit Reserve ab. \\
Transportauftrag & Zunaechst ueber Web-App & Externe Produktionsintegration bleibt spaeter moeglich. \\
Grundbedienung & Status, betrieblicher Stopp, Rueckkehr zur Ladestation & Stopp ist vom sicherheitsgerichteten Not-Halt zu unterscheiden. \\
\bottomrule
\end{longtable}

\section{Bedienung und Auftragserteilung}

Die erste Version verwendet eine Web-App als Bedien- und Auftragsoberflaeche. Ein Transportauftrag enthaelt im MVP mindestens Start, Ziel und Ausloesen des Auftrags. Spaeter koennen Auftrags-ID, Werkstuecktraeger-ID, Prioritaet oder weitere Prozessdaten hinzukommen.

Eine wichtige begriffliche Trennung wird bereits in Iteration A festgehalten:
\begin{quote}
Ein betrieblicher Stoppbefehl in einer Web-App ist kein sicherheitsgerichteter Not-Halt.
\end{quote}

Das Modell unterscheidet deshalb:
\begin{itemize}
  \item \emph{OperationalStop}: betrieblicher Stopp oder Pausenbefehl,
  \item \emph{SafetyStop}: sicherheitsbezogene Reaktion auf erkannte Gefahr oder unsichere Situation,
  \item \emph{EmergencyStop}: Not-Halt-Funktion, deren Bedarf und Ausfuehrung in der Sicherheitsbetrachtung festgelegt werden.
\end{itemize}

\section{Ladestation}

Die Funktion \enquote{Rueckkehr zur Ladestation} gehoert zum Zielsystem. Offen bleibt fuer den MVP, ob der Roboter nur zur Ladestation faehrt oder ob automatisches Docking und automatisches Laden bereits realisiert werden. Die Ladestation wird daher in Baseline A.1 als externes System mit relevanter Schnittstelle gefuehrt.

\section{Zentrale externe Schnittstellen}
\label{sec:interfaces}

\begin{longtable}{L{0.14\textwidth} L{0.22\textwidth} L{0.24\textwidth} L{0.26\textwidth}}
\toprule
ID & Schnittstelle & Richtung & Gegenstand \\
\midrule
IF-001 & Benutzer -- Web-App & bidirektional & Transportauftrag, Betriebsbefehl, Statusanzeige. \\
IF-002 & Web-App -- Roboter & bidirektional & Auftrag, Ziel, Status, Fehler, Grundbefehle. \\
IF-003 & Werkstuecktraeger -- Roboter & mechanisch & Auflage, Halten, Lage, Masse, Schwerpunkt. \\
IF-004 & Transfersystem -- Roboter & mechanisch/organisatorisch & Perspektivische Uebergabe eines Werkstuecktraegers. \\
IF-005 & Roboter -- Fahrumgebung & physisch/informatorisch & Bewegung, Wahrnehmung, Karte, Hindernisse, Tueren. \\
IF-006 & Roboter -- Personen & sicherheitsrelevant & Erkennung, Abstand, Bewegung, Warnung, sicherer Zustand. \\
IF-007 & Roboter -- Ladestation & physisch/energetisch & Rueckkehr, Positionierung, Energieuebergabe, Ladezustand. \\
\bottomrule
\end{longtable}

\section{Entscheidungen der Baseline A.1}

\begin{longtable}{L{0.14\textwidth} L{0.50\textwidth} L{0.20\textwidth}}
\toprule
ID & Entscheidung & Status \\
\midrule
\modeltable{decisions_baseline_summary}
\bottomrule
\end{longtable}

\section{Annahmen}

\begin{longtable}{L{0.14\textwidth} L{0.56\textwidth} L{0.16\textwidth}}
\toprule
ID & Annahme & Status \\
\midrule
\modeltable{assumptions_a1}
\bottomrule
\end{longtable}

\section{Nicht-Ziele}

Fuer die erste Entwicklungsstufe sind keine primaeren Ziele:
\begin{itemize}
  \item Ausseneinsatz,
  \item Treppensteigen,
  \item autonomer Personenverkehr,
  \item Manipulatorarm,
  \item humanoide Manipulation,
  \item vollstaendig autonome Be- und Entladung beliebiger Gegenstaende,
  \item mehrere Stockwerke per Aufzug,
  \item autonome Flottenoptimierung,
  \item Betrieb auf oeffentlichen Strassen,
  \item hochdynamische Fahrgeschwindigkeiten.
\end{itemize}

\section{Zentrale Risiken und offene Herausforderungen}

\begin{longtable}{L{0.14\textwidth} L{0.30\textwidth} L{0.40\textwidth}}
\toprule
ID & Herausforderung & Leitfrage \\
\midrule
R-01 & Navigation in veraenderlicher Umgebung & Wie zuverlaessig lassen sich Grundkarte und aktuelle Sensordaten kombinieren? \\
R-02 & Lokalisierung & Was passiert, wenn die reale Umgebung stark von der Karte abweicht? \\
R-03 & Hindernisse & Welche Hindernisse muss das System zuverlaessig erkennen? \\
R-04 & Personen im Fahrbereich & Wie wird sicherer gemeinsamer Betrieb mit Menschen erreicht? \\
R-05 & Materialtransport & Wie wird der Werkstuecktraeger mechanisch aufgenommen, gehalten und gegen Verrutschen gesichert? \\
R-06 & Blockierter Weg & Wann wartet der Roboter, wann plant er neu und wann bricht er ab? \\
R-07 & Energie & Wie wird die Betriebsdauer von ca. 2 h unter realistischem Nutzungsprofil nachgewiesen? \\
R-08 & Safety & Welche Sicherheitsarchitektur ist fuer realen Unterrichts- und Demonstrationsbetrieb erforderlich? \\
R-09 & Karte & Wann ist eine erkannte Aenderung temporaer und wann dauerhaft? \\
R-10 & Integration & Wie bleiben \sysml-Modell, Software, Elektronik, Mechanik und Tests synchron? \\
\bottomrule
\end{longtable}

\section{Definition of Done fuer Iteration A}

Iteration A ist mit Baseline A.1 fachlich geschlossen, weil Problem, Zweck, System of Interest, Systemgrenze, Einsatzumgebung, Stakeholder, Stakeholder Needs, Betriebsfaelle, Randbedingungen, Annahmen, Nicht-Ziele, Schnittstellen, Risiken und Entwicklungsprinzipien hinreichend beschrieben sind. Die naechste Iteration leitet daraus Use Cases, testbare Requirements und den formalen MVP-Schnitt ab.

\chapter{Stakeholder und Beduerfnisse}

\section{Stakeholderanalyse}

Stakeholder sind nicht nur direkte Nutzer. Fuer einen mobilen Roboter sind auch Betreiber, Wartung, Sicherheitsverantwortliche, IT, Entwicklung, Einkauf und Personen im Fahrbereich relevant. Das Lehrbuch hebt hervor, dass ein technisch funktionierendes System im Betrieb scheitern kann, wenn Wartbarkeit, Sicherheit, Datenschutz oder Betriebskontext nicht frueh beruecksichtigt werden \cite[Kap.~7]{becker2026sysml}.

\begin{longtable}{L{0.18\textwidth} L{0.30\textwidth} L{0.38\textwidth}}
\toprule
ID & Stakeholder & Interesse \\
\midrule
\modeltable{stakeholders_a1}
\bottomrule
\end{longtable}

\section{Stakeholder Needs}

\begin{longtable}{L{0.15\textwidth} L{0.18\textwidth} L{0.47\textwidth} L{0.10\textwidth}}
\toprule
ID & Quelle & Need & Prioritaet \\
\midrule
\modeltable{needs_a1_initial}
\bottomrule
\end{longtable}

\section{Welche Needs haben schon ein Requirement?}

Diese Tabelle ist keine Pflegeliste, sondern eine Antwort des Modells. Stakeholder sind dort Rollen (\texttt{part def}), Needs sind \texttt{concern}-Elemente mit ihren Stakeholdern, und jedes Requirement ist per Abhaengigkeit mit den Needs verbunden, aus denen es abgeleitet ist. Steht in der letzten Spalte \enquote{noch keins}, hat das Modell zu diesem Need noch keine Anforderung.

\begin{longtable}{L{0.18\textwidth} L{0.26\textwidth} L{0.42\textwidth}}
\toprule
Need & Stakeholder & Abgeleitete Requirements \\
\midrule
\modeltable{needs_coverage_a1}
\bottomrule
\end{longtable}

\section{Konflikte und Trade-offs}

\begin{longtable}{L{0.20\textwidth} L{0.32\textwidth} L{0.34\textwidth}}
\toprule
Konflikt & Beschreibung & Zu klaerende Entscheidung \\
\midrule
Sicherheit vs. Durchsatz & Niedrige Geschwindigkeit und konservatives Stoppen erhoehen Sicherheit, koennen aber Transportzeit verlaengern. & Zielgeschwindigkeit, Sicherheitsabstand, Zonenlogik \\
Kosten vs. Wahrnehmungsleistung & Hoeherwertige Sensorik verbessert Navigation und Hinderniserkennung, erhoeht aber Kosten und Energiebedarf. & Sensorvariante fuer MVP und Zielsystem \\
MVP-Fokus vs. Zielsystem & Ein kleiner MVP reduziert Aufwand, darf aber die spaetere Transportfunktion nicht architektonisch blockieren. & Welche Zielsystem-Schnittstellen werden bereits modelliert? \\
Kartierung vs. Stabilitaet & Automatische Kartenupdates koennen hilfreich sein, duerfen aber keine fehlerhaften dauerhaften Kartenveraenderungen erzeugen. & Trennung von globaler Karte, lokaler Karte und temporaeren Hindernissen \\
\bottomrule
\end{longtable}

\section{Detaillierte Stakeholder Needs aus Iteration A}

Die folgenden Needs wurden in der Vervollstaendigung von Iteration A feiner strukturiert. Sie ergaenzen die initiale Need-Liste und dienen als direkte Quelle fuer Iteration B.

\subsection{Bedienung und Auftrag}

\begin{longtable}{L{0.18\textwidth} L{0.16\textwidth} L{0.50\textwidth}}
\toprule
ID & Quelle & Need \\
\midrule
\modeltable{needs_a1_bedienung}
\bottomrule
\end{longtable}

\subsection{Transport}

\begin{longtable}{L{0.18\textwidth} L{0.16\textwidth} L{0.50\textwidth}}
\toprule
ID & Quelle & Need \\
\midrule
\modeltable{needs_a1_transport}
\bottomrule
\end{longtable}

\subsection{Navigation und Umgebung}

\begin{longtable}{L{0.18\textwidth} L{0.16\textwidth} L{0.50\textwidth}}
\toprule
ID & Quelle & Need \\
\midrule
\modeltable{needs_a1_navigation}
\bottomrule
\end{longtable}

\subsection{Safety}

\begin{longtable}{L{0.18\textwidth} L{0.16\textwidth} L{0.50\textwidth}}
\toprule
ID & Quelle & Need \\
\midrule
\modeltable{needs_a1_safety}
\bottomrule
\end{longtable}

\subsection{Betrieb, Wartung und Entwicklung}

\begin{longtable}{L{0.18\textwidth} L{0.16\textwidth} L{0.50\textwidth}}
\toprule
ID & Quelle & Need \\
\midrule
\modeltable{needs_a1_betrieb}
\bottomrule
\end{longtable}

\subsection{Unterricht und Demonstration}

\begin{longtable}{L{0.18\textwidth} L{0.16\textwidth} L{0.50\textwidth}}
\toprule
ID & Quelle & Need \\
\midrule
\modeltable{needs_a1_unterricht}
\bottomrule
\end{longtable}

\chapter{Requirements}

\section{Anforderungsprinzip}

Requirements beschreiben, was das System leisten oder einhalten muss. Technische Loesungen werden erst spaeter entschieden. Eine Aussage wie \enquote{Der Roboter verwendet einen bestimmten LiDAR} ist daher zunaechst keine Anforderung, sondern eine moegliche Architekturentscheidung. Diese Trennung entspricht der MBSE-Methodik des Lehrbuchs \cite[Kap.~6]{becker2026sysml}.

Jede Anforderung erhaelt:
\begin{itemize}
  \item eine stabile ID,
  \item einen kurzen Namen,
  \item einen Typ,
  \item eine Quelle oder ein Stakeholder Need,
  \item ein Akzeptanz- oder Nachweiskriterium,
  \item einen Status.
\end{itemize}

\section{Initiale Systemanforderungen}

\begin{longtable}{L{0.15\textwidth} L{0.16\textwidth} L{0.36\textwidth} L{0.18\textwidth}}
\toprule
ID & Typ & Requirement & Akzeptanzidee \\
\midrule
\modeltable{requirements_a1_initial}
\bottomrule
\end{longtable}

\section{MVP-relevante Anforderungen}

Fuer den ersten MVP werden zunaechst die Anforderungen priorisiert, die den technischen Kern aus Lokalisierung, Navigation, Umgebungserfassung, Hindernisbehandlung, Bewegungsausfuehrung und Statusmeldung pruefen.

\begin{longtable}{L{0.18\textwidth} L{0.16\textwidth} L{0.48\textwidth}}
\toprule
Requirement & MVP-Status & Begruendung \\
\midrule
\modeltable{mvp_a1}
\bottomrule
\end{longtable}

\section{Offene Quantifizierungen}

Folgende Anforderungen sind bewusst noch nicht numerisch festgelegt:
\begin{itemize}
  \item genaue Geschwindigkeitstoleranz und Fahrprofil fuer $0{,}5\,\mathrm{m/s}$,
  \item erlaubte Positionsabweichung,
  \item minimale Hindernisdetektionsdistanz,
  \item Bremsweg und Sicherheitsabstand,
  \item Nutzlast und Kistengeometrie,
  \item Batterieauslegung fuer das 2-h-Unterrichtsprofil,
  \item maximale Abmessungen des Roboters,
  \item Verfuegbarkeit und zulassige Fehlerrate.
\end{itemize}

Diese Werte werden nicht geraten, sondern in einer spaeteren Iteration aus Nutzungskontext, Safety-Konzept, mechanischem Aufbau und Testmoeglichkeiten abgeleitet.

\section{Ergaenzte Anforderungen aus Iteration A}

Aus den detaillierten Needs, Betriebsfaellen und dem Kartenkonzept ergeben sich folgende zusaetzliche oder praezisierte Anforderungen fuer die naechste Arbeitsfassung:

\begin{longtable}{L{0.15\textwidth} L{0.16\textwidth} L{0.36\textwidth} L{0.18\textwidth}}
\toprule
ID & Typ & Requirement & Akzeptanzidee \\
\midrule
\modeltable{requirements_a1_iteration_a}
\bottomrule
\end{longtable}

\section{Randbedingungen als Requirements-Quellen}

\begin{longtable}{L{0.15\textwidth} L{0.35\textwidth} L{0.32\textwidth}}
\toprule
Randbedingung & Bedeutung & Betroffene Requirements \\
\midrule
C-001 & Indoor-Betrieb & REQ-NAV-001, REQ-SAFE-001, REQ-MAP-001 \\
C-002 & Autonomer Fahrbetrieb & REQ-MIS-001, REQ-NAV-001, REQ-SAFE-001 \\
C-003 & Definierte Ziele & REQ-MIS-001, REQ-NAV-001 \\
C-004 & Vorhandene Grundkarte & REQ-MAP-001, REQ-LOC-001 \\
C-005 & Dynamische Umgebung & REQ-MAP-002, REQ-MAP-003, REQ-NAV-002, REQ-SAFE-001 \\
\bottomrule
\end{longtable}

\chapter{Betriebsfaelle und Use Cases}

\section{Use-Case-Prinzip}

Use Cases beschreiben Ziele und Interaktionen aus Sicht eines Akteurs. Sie sind keine Algorithmen und ersetzen keine Architektur. Das Lehrbuch betont, dass Use Cases helfen, Systemgrenzen, Akteure, Vorbedingungen, Normalablaeufe, Ausnahmen und erwartete Ergebnisse sichtbar zu machen \cite[Kap.~7]{becker2026sysml}.

\section{Betriebsfaelle aus Baseline A.1}

In Iteration A werden zuerst reale Betriebssituationen gesammelt. Diese werden in Iteration B zu vollstaendigen Use Cases und anschliessend zu Requirements verfeinert.

\begin{figure}[htbp]
  \centering
  \resizebox{\textwidth}{!}{\begin{tikzpicture}[
  op/.style={draw, rounded corners=2pt, align=center, minimum width=24mm, minimum height=8mm, fill=gray!8, font=\small},
  alt/.style={draw, rounded corners=2pt, align=center, minimum width=26mm, minimum height=8mm, fill=yellow!12, font=\small},
  arrow/.style={-{Latex[length=2mm]}, thick}
]
\node[op] (load) {Werkstueck-\\traeger aufsetzen};
\node[op, right=7mm of load] (order) {Auftrag\\Web-App};
\node[op, right=7mm of order] (localize) {Karte laden\\lokalisieren};
\node[op, right=7mm of localize] (drive) {autonom\\fahren};
\node[op, right=7mm of drive] (door) {offene Tuer\\Flur};
\node[op, right=7mm of door] (target) {Ziel erreichen\\abschliessen};

\node[alt, below=12mm of drive] (obstacle) {Hindernis\\erkannt};
\node[alt, below=12mm of obstacle] (blocked) {umfahren,\\neu planen\\oder stoppen};
\node[alt, below=12mm of target] (charge) {Rueckkehr\\Ladestation};

\draw[arrow] (load) -- (order);
\draw[arrow] (order) -- (localize);
\draw[arrow] (localize) -- (drive);
\draw[arrow] (drive) -- (door);
\draw[arrow] (door) -- (target);
\draw[arrow] (drive) -- (obstacle);
\draw[arrow] (obstacle) -- (blocked);
\draw[arrow] (blocked.east) -| (target.south);
\draw[arrow] (target) -- (charge);
\end{tikzpicture}
}
  \caption{End-to-End-Betriebsablauf des ersten Demonstrators}
  \label{fig:operating-cases}
\end{figure}

\begin{longtable}{L{0.14\textwidth} L{0.42\textwidth} L{0.30\textwidth}}
\toprule
ID & Betriebsfall & Rolle fuer Iteration B \\
\midrule
OP-01 & Transportauftrag ueber Web-App erteilen & Quelle fuer Auftrags- und UI-Requirements. \\
OP-02 & Werkstuecktraeger manuell aufsetzen & Quelle fuer mechanische Schnittstelle und Bedienprozess. \\
OP-03 & Werkstuecktraeger durch externes Transfersystem uebergeben & Future-Schnittstelle fuer Transfersystemintegration. \\
OP-04 & Autonome Fahrt zu einem Ziel & Zentraler Navigations-Use-Case. \\
OP-05 & Raum verlassen und offene Tuer durchfahren & Randbedingung fuer Navigation, Abmessungen und Tests. \\
OP-06 & Flur befahren & Betriebsumgebung mit Personen- und Hindernisrisiko. \\
OP-07 & Anderen Raum anfahren & End-to-End-Nachweis ueber mehrere Bereiche. \\
OP-08 & Temporaeres Hindernis erkennen & Quelle fuer Perception- und Kartenanforderungen. \\
OP-09 & Hindernis umfahren & Quelle fuer lokale Planung und Safety-Abgrenzung. \\
OP-10 & Blockierten Fahrweg behandeln & Quelle fuer Abbruch-, Warte- und Statuslogik. \\
OP-11 & Person im Fahrbereich behandeln & Quelle fuer Safety- und Betriebsanforderungen. \\
OP-12 & Ziel erreichen und Auftrag abschliessen & Quelle fuer Missionsstatus und Abnahmekriterium. \\
OP-13 & Benutzer fordert betrieblichen Stopp & Trennung OperationalStop, SafetyStop und EmergencyStop. \\
OP-14 & Benutzer fordert Rueckkehr zur Ladestation & Quelle fuer Energie- und Ladestationsschnittstelle. \\
OP-15 & Technischer Fehler tritt auf & Quelle fuer Diagnose und sicheren Zustand. \\
OP-16 & Roboter muss sicheren Zustand erreichen & Quelle fuer Safety-Architektur und Verifikation. \\
\bottomrule
\end{longtable}

\section{Use-Case-Uebersicht fuer Iteration B}

\begin{longtable}{L{0.15\textwidth} L{0.32\textwidth} L{0.34\textwidth} L{0.10\textwidth}}
\toprule
ID & Use Case & Ziel & MVP \\
\midrule
UC-001 & Transportauftrag erteilen & Bediener gibt Start und Ziel ueber die Web-App vor. & Ja \\
UC-002 & Werkstuecktraeger bereitstellen & Mensch oder Transfersystem platziert den Werkstuecktraeger auf dem Roboter. & Teilweise \\
UC-003 & Transportziel anfahren & Roboter faehrt autonom zu einem definierten Ziel. & Ja \\
UC-004 & Raeume und Flur durchfahren & Roboter bewegt sich durch offene Tueren und Flurabschnitte. & Ja \\
UC-005 & Grundkarte laden & Roboter laedt eine vorhandene Karte des Betriebsbereichs. & Ja \\
UC-006 & Position bestimmen & Roboter lokalisiert sich in der Karte. & Ja \\
UC-007 & Hindernis behandeln & Roboter erkennt Hindernisse und reagiert durch Stoppen, Umfahren oder Neuplanung. & Ja \\
UC-008 & Blockierten Fahrweg behandeln & Roboter erkennt, dass keine sinnvolle Umfahrung moeglich ist, und meldet Blockade. & Ja \\
UC-009 & Betrieblichen Stopp ausloesen & Benutzer stoppt den laufenden Betrieb kontrolliert. & Ja \\
UC-010 & Rueckkehr zur Ladestation ausloesen & Roboter faehrt auf Benutzerbefehl oder spaeter automatisch zur Ladestation. & Teilweise \\
UC-011 & Batterie laden & Roboter laedt autonom oder assistiert. & Spaeter \\
UC-012 & Fehler diagnostizieren & Wartung erhaelt Informationen zur Fehlersuche. & Teilweise \\
UC-013 & Karte verwalten & Grundkarte, lokale Hindernisse und dauerhafte Aenderungen bleiben kontrollierbar. & Teilweise \\
\bottomrule
\end{longtable}

\section{Ausgearbeiteter Use Case UC-003: Transportziel anfahren}

\begin{longtable}{L{0.24\textwidth} L{0.64\textwidth}}
\toprule
Feld & Beschreibung \\
\midrule
ID & UC-003 \\
Name & Transportziel anfahren \\
Primaerer Akteur & Bediener oder Auftragsquelle \\
Ziel & Roboter erreicht einen definierten Zielstandort autonom. \\
Vorbedingungen & Karte ist geladen; Roboter ist lokalisiert; Zielstandort ist bekannt; System befindet sich in betriebsbereitem Zustand. \\
Normalablauf & Auftrag wird angenommen; Ziel wird validiert; globale Route wird geplant; lokale Bewegung wird fortlaufend geplant; Roboter bewegt sich zum Ziel; Status wird aktualisiert; Zielerreichung wird gemeldet. \\
Alternativen & Hindernis wird erkannt und umfahren; Ziel ist nicht erreichbar; Lokalisierung ist unsicher; Fahrweg ist blockiert. \\
Ergebnis & Ziel erreicht oder begruendete Nichtausfuehrung mit Statusmeldung. \\
Verknuepfte Requirements & REQ-MIS-001, REQ-MAP-001, REQ-LOC-001, REQ-NAV-001, REQ-NAV-002, REQ-STA-001 \\
\bottomrule
\end{longtable}

\section{Ausgearbeiteter Use Case UC-007: Hindernis behandeln}

\begin{longtable}{L{0.24\textwidth} L{0.64\textwidth}}
\toprule
Feld & Beschreibung \\
\midrule
ID & UC-007 \\
Name & Hindernis behandeln \\
Primaerer Akteur & Umgebung als externes Ereignis; indirekt Personen im Fahrbereich \\
Ziel & Der Roboter vermeidet Kollisionen und setzt die Mission fort, falls dies sicher moeglich ist. \\
Vorbedingungen & Roboter befindet sich in Fahrt oder geplanter Bewegung; Sensorik liefert Umgebungsdaten. \\
Normalablauf & Hindernis wird erkannt; Hindernis wird in lokaler Karte beruecksichtigt; lokale Bewegung wird angepasst; Mission laeuft weiter. \\
Alternativen & Keine sichere Umfahrung moeglich; Hindernis befindet sich zu nah; Sensorik unsicher; Roboter stoppt und meldet Blockade. \\
Ergebnis & Umfahrung oder sicherer Stopp mit Statusmeldung. \\
Verknuepfte Requirements & REQ-MAP-002, REQ-NAV-002, REQ-SAFE-001, REQ-STA-001 \\
\bottomrule
\end{longtable}

\chapter{MVP-Abgrenzung}

\section{MVP-Grundsatz}

Der MVP ist der kleinste durchgaengige Systemausschnitt, mit dem die zentrale technische Hypothese geprueft werden kann:
\begin{quote}
Ein mobiler Indoor-Roboter kann einen manuell aufgesetzten Werkstuecktraeger als Payload mitfuehren, ueber eine Web-App einen Transportauftrag erhalten, auf Basis einer vorhandenen Grundkarte lokalisiert werden, ein Ziel autonom anfahren, offene Tueren und einen Flur durchfahren, Veraenderungen der Umgebung wahrnehmen, Hindernisse behandeln und seinen Missionsstatus nachvollziehbar melden.
\end{quote}

Der MVP ist damit kein separates Produkt und kein Wegwerfprototyp. Er nutzt dieselbe Zielarchitektur, implementiert aber nur einen bewusst begrenzten Umfang.

\section{MVP 1: Navigationskern}

\begin{longtable}{L{0.42\textwidth} L{0.42\textwidth}}
\toprule
Im MVP enthalten & Im MVP nicht enthalten \\
\midrule
Fahren auf ebener Indoor-Flaeche & Vollautomatische Kistenaufnahme \\
Grundkarte laden & Vollautomatische Kistenuebergabe \\
Lokalisierung in der Karte & Autonomes Docking und Laden \\
Zielstandort ueber Web-App empfangen & Flottenmanagement \\
Autonomes Zielanfahren & Aufzug- oder Tuerschnittstellen \\
Offene Tueren und Flur durchfahren & Automatisches Oeffnen geschlossener Tueren \\
Hindernisse erkennen & Dynamische Auftragsoptimierung \\
Stoppen oder lokale Umfahrung & Personenerkennung per KI als separate Klassifikation \\
Missionsstatus melden & Zertifizierte Safety-Architektur \\
Rueckkehr zur Ladestation als Fahrziel & Vollautomatisches Laden als Muss-Funktion \\
Minimale Diagnose fuer Tests & Vollstaendige Wartungs- und Betriebsoberflaeche \\
\bottomrule
\end{longtable}

\section{MVP-Erfolgskriterien}

\begin{longtable}{L{0.14\textwidth} L{0.48\textwidth} L{0.26\textwidth}}
\toprule
ID & Erfolgskriterium & Nachweismethode \\
\midrule
MVP-SC-001 & Roboter laedt eine bereitgestellte Karte und kann sich in ihr lokalisieren. & Integrationstest \\
MVP-SC-002 & Roboter faehrt autonom zu mindestens einem definierten Ziel. & Systemtest \\
MVP-SC-003 & Roboter erkennt ein Hindernis im Fahrweg und vermeidet Kollision durch Stoppen oder Umfahren. & Systemtest oder Simulation \\
MVP-SC-004 & Roboter meldet Missionsstatus: bereit, unterwegs, Ziel erreicht, blockiert, Fehler. & Integrationstest \\
MVP-SC-005 & Zu allen MVP-Anforderungen existiert eine Traceability-Kette bis zu Testfaellen. & Review \\
MVP-SC-006 & Roboter durchfaehrt im Testaufbau mindestens eine offene Tuer und einen Flurabschnitt. & Systemtest \\
MVP-SC-007 & Roboter erreicht fuer ein definiertes Unterrichtsprofil eine Betriebsdauer von ca. 2 h. & Systemtest \\
\bottomrule
\end{longtable}

\section{Begruendung der Abgrenzung}

Der Kern des Zielsystems liegt in der Kopplung von Karte, Lokalisierung, Umgebungserfassung, Bewegungsplanung, Bewegungsausfuehrung und Statusmeldung. Der physische Materialumschlag ist wichtig, wuerde jedoch im ersten Schritt zusaetzliche mechanische, elektrische und sicherheitstechnische Komplexitaet einbringen. Die MVP-Abgrenzung reduziert diese Komplexitaet, ohne die Zielarchitektur zu verwerfen.

\section{Risiko der MVP-Abgrenzung}

Das groesste Risiko besteht darin, dass die spaetere Transportaufnahme zu spaet beruecksichtigt wird. Deshalb bleibt das \emph{TransportSystem} bereits Teil der logischen Zielarchitektur. Es wird im MVP noch nicht voll implementiert, aber seine Schnittstellen zu Mission, Navigation, Energie und Safety werden frueh sichtbar gemacht.

\chapter{Systemarchitektur}

\section{Architekturprinzip}

Die Architektur wird zunaechst funktional und logisch beschrieben. Konkrete Komponenten, Sensoren, Rechner oder \ros-Pakete werden erst danach ausgewaehlt. Dadurch bleibt das Modell loesungsneutral genug, um Alternativen bewerten zu koennen, aber konkret genug, um Schnittstellen und Tests vorzubereiten.

\section{Hauptfunktionen}

\begin{longtable}{L{0.16\textwidth} L{0.30\textwidth} L{0.38\textwidth}}
\toprule
ID & Funktion & Zweck \\
\midrule
FUN-MIS-001 & Transportauftrag bearbeiten & Auftrag annehmen, pruefen, ausfuehren und Status verwalten \\
FUN-MAP-001 & Grundkarte verwalten & Karte laden, referenzieren und fuer Lokalisierung bereitstellen \\
FUN-LOC-001 & Position bestimmen & Roboterpose relativ zur Karte schaetzen \\
FUN-PER-001 & Umgebung erfassen & Sensorinformationen fuer Navigation und Hindernisbehandlung bereitstellen \\
FUN-NAV-001 & Route und Bewegung planen & Globales Ziel und lokale Hindernisse in Bewegungsbefehle ueberfuehren \\
FUN-MOT-001 & Bewegung ausfuehren & Bewegungsbefehle in Motorsteuerung und Radbewegung umsetzen \\
FUN-SAFE-001 & Sicheren Zustand herstellen & Bei nicht beherrschbaren Situationen stoppen oder sicheren Zustand melden \\
FUN-STA-001 & Status und Diagnose melden & Betriebs- und Missionsinformationen fuer Nutzer, Tests und Wartung bereitstellen \\
FUN-TRN-001 & Materialkiste handhaben & Kiste aufnehmen, sichern und abgeben \\
\bottomrule
\end{longtable}

\section{Logische Teilsysteme}

\begin{longtable}{L{0.22\textwidth} L{0.46\textwidth} L{0.18\textwidth}}
\toprule
Teilsystem & Verantwortung & MVP-Status \\
\midrule
MissionSystem & Auftragsannahme, Missionszustand, Zielverwaltung, Ablaufkoordination & Teilweise \\
TransportSystem & Mechanische und logische Handhabung der Materialkiste & Zielsystem \\
NavigationSystem & Lokalisierung, globale Planung, lokale Planung, Zielerreichung & MVP \\
PerceptionSystem & Sensoranbindung, Hinderniserkennung, lokale Umgebungskarte & MVP \\
MotionSystem & Bewegungsregelung, Motoransteuerung, Odometriedaten & MVP \\
EnergySystem & Batterieueberwachung, Energiezustand, spaeter Laden/Docking & Minimal \\
SafetySystem & Ueberwachung sicherheitsrelevanter Zustaende, sicherer Stopp & MVP reduziert \\
CommunicationSystem & Status, Diagnose, Auftrags- und Datenschnittstellen & Minimal \\
UserInterface & Bedien- und Statusschnittstelle fuer Mensch oder externes System & Minimal \\
\bottomrule
\end{longtable}

\section{Architekturentscheidung ADR-0001}

\begin{decisionbox}{Entscheidung}
Das Gesamtmodell wird von Beginn an mit allen zentralen Zielsystem-Teilsystemen aufgebaut. Der MVP implementiert nur einen markierten Teil davon.
\end{decisionbox}

\begin{longtable}{L{0.20\textwidth} L{0.64\textwidth}}
\toprule
Aspekt & Inhalt \\
\midrule
Status & Angenommen fuer Version 0.1 \\
Alternativen & Separates MVP-Modell; direkte Implementierung ohne Zielarchitektur; vollstaendige Zielarchitektur vor erstem Prototyp \\
Begruendung & Ein separates MVP-Modell wuerde spaeter Doppelarbeit erzeugen. Eine direkte Implementierung wuerde Traceability verlieren. Eine vollstaendige Zielarchitektur vor dem ersten Prototyp waere zu schwerfaellig. \\
Auswirkung & Alle Artefakte benoetigen Status- und Prioritaetsattribute. MVP-Elemente bleiben auf Zielsystem-Elemente rueckfuehrbar. \\
Risiko & Noch nicht implementierte Zielsystembereiche koennen zu grob bleiben. Dieses Risiko wird durch offene Punkte und spaetere Iterationen kontrolliert. \\
\bottomrule
\end{longtable}

\section{Erste Schnittstellen}

\begin{longtable}{L{0.16\textwidth} L{0.20\textwidth} L{0.20\textwidth} L{0.28\textwidth}}
\toprule
ID & Von & Nach & Inhalt \\
\midrule
IF-MIS-NAV-001 & MissionSystem & NavigationSystem & Zielstandort, Missionsziel, Abbruchsignal \\
IF-NAV-MOT-001 & NavigationSystem & MotionSystem & Bewegungsbefehl, Geschwindigkeitsvorgabe, Stoppanforderung \\
IF-MOT-NAV-001 & MotionSystem & NavigationSystem & Odometriedaten, Bewegungsstatus \\
IF-PER-NAV-001 & PerceptionSystem & NavigationSystem & Hindernisse, lokale Karte, freie Flaechen \\
IF-SAFE-MOT-001 & SafetySystem & MotionSystem & sicherheitsrelevanter Stopp oder Freigabe \\
IF-ENG-MIS-001 & EnergySystem & MissionSystem & Energiezustand, Warnung niedriger Ladezustand \\
IF-MIS-UI-001 & MissionSystem & UserInterface & Missionsstatus, Fehler, Zielerreichung \\
\bottomrule
\end{longtable}

\chapter{Hardware-/Software-Allokation und ROS2-Mapping}

\section{Allokationsprinzip}

Die Hardware-/Software-Allokation wird aus Funktionen abgeleitet. Erst wenn die Funktion, ihr Nachweis und ihre Schnittstellen verstanden sind, werden konkrete Komponenten festgelegt. Dies ist besonders wichtig fuer Sensorik, Rechnerplattform, Motorregelung und Safety.

\section{Vorlaeufige Allokation}

\begin{longtable}{L{0.18\textwidth} L{0.26\textwidth} L{0.28\textwidth} L{0.14\textwidth}}
\toprule
Funktion & Hardware-Anteil & Software-Anteil & Status \\
\midrule
FUN-MAP-001 & Onboard-Speicher, Rechner & Kartenserver, Kartenverwaltung, Versionierung & MVP \\
FUN-LOC-001 & Sensorik, Encoder, IMU, Rechner & Lokalisierungsalgorithmus, Transformationsbaum & MVP \\
FUN-PER-001 & LiDAR/Kamera/weitere Sensorik & Sensortreiber, Filterung, Hindernisextraktion & MVP \\
FUN-NAV-001 & Rechner & globale Planung, lokale Planung, Recovery-Verhalten & MVP \\
FUN-MOT-001 & Motoren, Motorcontroller, Encoder & hardwarenahe Regelung, \ros-Anbindung & MVP \\
FUN-SAFE-001 & Sensorik, Not-Halt, ggf. Safety-Hardware & Ueberwachung, Stopp-Logik, Fehlerstatus & MVP reduziert \\
FUN-TRN-001 & Mechanik der Aufnahme, Aktoren, Sensorik & Handhabungslogik, Verriegelungsstatus & Zielsystem \\
\bottomrule
\end{longtable}

\section{ROS2-Mapping}

Das Lehrbuch beschreibt das Mapping von \sysml-Elementen auf \ros-Elemente als Architekturentscheidung, nicht als starre Uebersetzungsregel \cite[Kap.~16]{becker2026sysml}. Ein \sysml-Element kann mehrere Nodes beeinflussen, und ein Node kann mehrere modellierte Verantwortlichkeiten buendeln.

\begin{longtable}{L{0.20\textwidth} L{0.20\textwidth} L{0.22\textwidth} L{0.22\textwidth}}
\toprule
\sysml-Element & \ros-Element & Schnittstellenidee & Verknuepfte Anforderungen \\
\midrule
MissionSystem & \texttt{mission\_manager\_node} & Auftrag, Ziel, Mission State & REQ-MIS-001, REQ-MIS-002, REQ-STA-001 \\
UserInterfaceSystem, CommunicationSystem & \texttt{web\_interface\_node} & Auftraege und Bedienbefehle von der Web-App, Status und Diagnose zur Web-App & REQ-UI-001, REQ-UI-002 \\
NavigationSystem & \texttt{navigation\_node}, z.\,B. auf Basis Nav2 & Zielaktion, Pfad, Bewegungsbefehl & REQ-NAV-001 bis REQ-NAV-005, REQ-PERF-001 \\
NavigationSystem (Kartenverwaltung) & \texttt{map\_management\_node} & Karte bereitstellen, Kartenupdate-Kandidaten & REQ-MAP-001, REQ-MAP-003 \\
NavigationSystem (Lokalisierung) & \texttt{localization\_node} & Pose, Transformationsdaten & REQ-LOC-001 \\
PerceptionSystem & \texttt{perception\_node}, ggf. aufgeteilt in Sensortreiber und Hindernisverarbeitung & Sensordaten, Hindernisse, lokale Karte & REQ-MAP-002, REQ-SAFE-002 \\
MotionSystem & \texttt{motion\_controller\_node}, ggf. mit MCU-Firmware & \texttt{cmd\_vel}, Odometrie, Motorstatus & REQ-NAV-001 \\
SafetySystem & \texttt{safety\_supervisor\_node}, spaeter plus Hardware-Stopp & Safety-Zustand, Stoppgrund & REQ-SAFE-001, REQ-SAFE-002 \\
EnergySystem & \texttt{energy\_monitor\_node} & Batterie- und Warnstatus & REQ-PWR-001, REQ-PWR-002 \\
DiagnosticsSystem & \texttt{diagnostics\_node} & Diagnosemeldungen & REQ-DIAG-001 \\
\bottomrule
\end{longtable}

\section{Vorlaeufige ROS2-Schnittstellen}

\begin{longtable}{L{0.18\textwidth} L{0.20\textwidth} L{0.16\textwidth} L{0.30\textwidth}}
\toprule
Name & Typ & Richtung & Zweck \\
\midrule
\texttt{/goal\_pose} & Topic/Action offen & Mission zu Navigation & Zielvorgabe fuer Navigation \\
\texttt{/cmd\_vel} & Topic & Navigation zu Motion & Geschwindigkeitsvorgabe \\
\texttt{/odom} & Topic & Motion zu Navigation & Odometriedaten \\
\texttt{/map} & Topic/Service & Kartenverwaltung zu Navigation/Lokalisierung & Grundkarte \\
\texttt{/local\_obstacles} & Topic & Perception zu Navigation & erkannte Hindernisse \\
\texttt{/mission\_state} & Topic & Mission zu UI/Logging & Missionsstatus \\
\texttt{/safety\_state} & Topic & Safety zu Mission/Motion & Safety-Zustand und Stoppgrund \\
\texttt{/battery\_state} & Topic & Energy zu Mission/UI & Energiezustand \\
\texttt{/transport\_order} & Topic & Web-Interface zu Mission & Transportauftrag \\
\texttt{/operator\_command} & Topic & Web-Interface zu Mission & Status, Stopp, Fortsetzung, Rueckkehr \\
\texttt{/navigation\_result} & Topic & Navigation zu Mission & Ziel erreicht, blockiert, nicht erreichbar \\
\texttt{/robot\_pose} & Topic & Lokalisierung zu Navigation & geschaetzte Pose \\
\texttt{/map\_update\_candidate} & Topic & Kartenverwaltung & vorgeschlagene dauerhafte Kartenaenderung \\
\texttt{/diagnostics} & Topic & Diagnose zu Web-Interface & Diagnosemeldungen \\
\bottomrule
\end{longtable}

\section{Noch nicht entschiedene Alternativen}

\begin{itemize}
  \item Rechnerplattform: Raspberry Pi, Jetson, x86 oder Kombination mit MCU.
  \item Sensorik: 2D-LiDAR, Tiefenkamera, Stereo, Ultraschall oder kombinierte Sensorik.
  \item Motorregelung: reine Motorcontroller-Loesung, MCU-Firmware oder integrierte Servoregler.
  \item Navigation: ROS2 Nav2 als Referenz, Eigenanteile fuer lokale Hindernislogik, alternative SLAM-/Lokalisierungspakete.
  \item Simulation: Gazebo, Isaac Sim oder reduzierter Testsimulator fuer fruehe MVP-Tests.
\end{itemize}

\chapter{Verifikation, Validierung und Traceability}

\section{Verifikationsprinzip}

Ein Systemmodell ist nur dann nuetzlich, wenn es bis zum Nachweis fuehrt. Verifikation prueft, ob das System die spezifizierten Anforderungen erfuellt. Validierung prueft, ob das entwickelte System das tatsaechliche Problem der Stakeholder loest. Diese Unterscheidung entspricht dem im Lehrbuch beschriebenen Verifikations- und Validierungsverstaendnis \cite[Kap.~19]{becker2026sysml}.

\section{Initiale Testfaelle}

\begin{longtable}{L{0.14\textwidth} L{0.24\textwidth} L{0.22\textwidth} L{0.24\textwidth}}
\toprule
ID & Testfall & Verifiziert & Erwartetes Ergebnis \\
\midrule
\modeltable{tests_a1}
\bottomrule
\end{longtable}

Alle Testfaelle sind als Verifikationsdefinitionen im \sysml-Modell (\texttt{06\_tests.sysml}) hinterlegt. Keiner ist bisher an einem realen Roboter oder in einem dokumentierten Simulationslauf ausgefuehrt worden.

\section{Simulationspruefstand und Modellpruefung}

Gemaess DEC-014 wird zunaechst ausschliesslich in der Simulation verifiziert. Das Modell enthaelt dafuer einen logischen Pruefstand (\texttt{SimulationTestBench}) aus Roboter, Simulationswelt, Steuerungsadapter und Beobachter. Simulator, Weltgeometrie, Sensor- und Energiemodell sind noch nicht ausgewaehlt.

TC-NAV-002 und TC-UI-003 besitzen auswertbare Kriterien. Ohne vollstaendige Nachweise lautet das Ergebnis \enquote{inconclusive}, nicht \enquote{pass}. Die Modellpruefung (\texttt{tools/validate\_sysml.py}) laedt das Gesamtmodell mit der offiziellen SysML-v2-Referenzimplementierung, wertet synthetische Testdaten aus und prueft absichtlich beschaedigte Modellkopien. Synthetische Testdaten sind keine Beobachtungen.

\section{Testfallstruktur}

Jeder Testfall wird spaeter in folgender Struktur beschrieben:
\begin{itemize}
  \item Test-ID,
  \item verifizierte Requirement-ID,
  \item Ziel,
  \item Vorbedingungen,
  \item Testaufbau oder Simulationsumgebung,
  \item Stimulus oder Ablauf,
  \item erwartetes Ergebnis,
  \item Bewertungskriterium,
  \item benoetigte Messdaten,
  \item Status und letzter Durchlauf.
\end{itemize}

\section{Traceability-Kette}

Die minimale Traceability-Kette fuer dieses Projekt lautet:
\begin{quote}
Stakeholder Need $\rightarrow$ Requirement $\rightarrow$ Use Case oder Funktion $\rightarrow$ logisches Systemelement $\rightarrow$ Schnittstelle oder \ros-Element $\rightarrow$ Testfall.
\end{quote}

\begin{figure}[htbp]
  \centering
  \resizebox{\textwidth}{!}{\begin{tikzpicture}[
  item/.style={draw, rounded corners=2pt, align=center, minimum width=26mm, minimum height=8mm, fill=gray!8},
  arrow/.style={-{Latex[length=2mm]}, thick}
]
\node[item] (need) {Stakeholder\\Need};
\node[item, right=7mm of need] (req) {Requirement};
\node[item, right=7mm of req] (uc) {Use Case\\Betriebsfall};
\node[item, right=7mm of uc] (fun) {Funktion\\Verhalten};
\node[item, right=7mm of fun] (arch) {Systemelement\\Schnittstelle};
\node[item, right=7mm of arch] (test) {Test\\Nachweis};

\draw[arrow] (need) -- (req);
\draw[arrow] (req) -- (uc);
\draw[arrow] (uc) -- (fun);
\draw[arrow] (fun) -- (arch);
\draw[arrow] (arch) -- (test);
\end{tikzpicture}
}
  \caption{Minimale Traceability-Kette der Dokumentation}
  \label{fig:traceability-chain}
\end{figure}

\section{Traceability-Regeln}

\begin{itemize}
  \item Keine MVP-Anforderung ohne Testfall.
  \item Kein \ros-Element ohne modellierte Verantwortung.
  \item Keine Schnittstelle ohne Producer, Consumer und Zweck.
  \item Keine Architekturentscheidung ohne betrachtete Alternative.
  \item Jede geaenderte Anforderung loest eine Impact-Analyse auf Funktionen, Schnittstellen, Tests und Dokumentation aus.
\end{itemize}

\chapter{Iteration A: Status und Baseline-Freigabe}

\section{Ziel der Iteration}

Iteration A klaert den Projektanfang. Sie liefert noch keine endgueltige technische Architektur, sondern eine belastbare Grundlage fuer Requirements, Use Cases und den spaeteren MVP-Schnitt. Der Fokus liegt auf Zweck, Kontext, Systemgrenze, Stakeholdern, Needs, Annahmen, Randbedingungen, Nicht-Zielen, Risiken, Schnittstellen und Entwicklungsprinzipien.

\section{Baseline A.1}

Der Stand dieser Dokumentation wird als \emph{System Definition Baseline A.1} gefuehrt. Die Baseline enthaelt die Festlegungen aus dem aktuellen Projektgespraech:
\begin{itemize}
  \item Werkstuecktraeger eines Transfersystems als Transportgut,
  \item externe Beladung durch Mensch oder Transfersystem,
  \item produktionsaehnlicher Schulbereich als Demonstratorumgebung,
  \item mehrere Raeume und ein Flur als Fahrbereich,
  \item Befahren offener Tueren,
  \item Nominalgeschwindigkeit des MVP von ca. $0{,}5\,\mathrm{m/s}$,
  \item Betriebsdauer der ersten Version von ca. 2 h,
  \item Auftragserteilung ueber Web-App,
  \item Grundbedienung mit Status, betrieblichem Stopp und Rueckkehr zur Ladestation.
\end{itemize}

\section{Ergebnisstand}

\begin{longtable}{L{0.32\textwidth} L{0.22\textwidth} L{0.32\textwidth}}
\toprule
Element & Stand & Bemerkung \\
\midrule
Problemstellung & Definiert & Materialtransport in semi-dynamischer Indoor-Umgebung. \\
Systemzweck & Definiert & Autonomer Transport von Werkstuecktraegern zwischen definierten Standorten. \\
System of Interest & Definiert & Mobiler Roboter mit notwendiger Onboard-Hardware und Onboard-Software. \\
Systemgrenze & Definiert fuer Baseline A.1 & Ladestation, Web-App, Transfersystem und Umgebung als externe Schnittstellen sichtbar. \\
Einsatzumgebung & Definiert & Produktionsaehnlicher Schulbereich, mehrere Raeume, Flur, offene Tueren. \\
Kartenkonzept & Definiert auf Systemebene & Trennung von Basisinformation, aktueller Umgebung und dauerhaften Aenderungen. \\
Hauptakteure & Definiert & Bediener, Empfaenger, Betreiber, Personen, Wartung, IT, Safety, Entwicklung, Einkauf, Unterricht. \\
Stakeholder Needs & Erster Satz vorhanden & Needs sind Quelle fuer Iteration B. \\
Betriebsfaelle & Definiert & OP-01 bis OP-16 als Grundlage fuer Use Cases. \\
Randbedingungen & Teilweise quantifiziert & Geschwindigkeit ca. $0{,}5\,\mathrm{m/s}$, Betriebsdauer ca. 2 h. \\
Annahmen & Dokumentiert & A-001 bis A-011. \\
Nicht-Ziele & Definiert & Verhindert fruehes Ausufern des Scopes. \\
Schnittstellen & Definiert auf Kontextebene & IF-001 bis IF-007. \\
Hauptrisiken & Identifiziert & R-01 bis R-10. \\
Entwicklungsprinzipien & Definiert & DP-01 bis DP-07. \\
Technische Architektur & Bewusst offen & Keine vorzeitige Festlegung auf Sensorik, Rechner, Fahrwerk oder \ros-Paket. \\
\bottomrule
\end{longtable}

\section{Freigabeentscheidung}

Iteration A ist fachlich geschlossen. Die noch offenen Punkte verhindern nicht den Uebergang in Iteration B, weil sie als Open Issues dokumentiert sind und dort in testbare Anforderungen ueberfuehrt oder durch Entscheidungen aufgeloest werden.

\section{Umgang mit Open Issues}

Ein Open Issue ist in dieser Dokumentation keine vergessene Luecke, sondern eine bewusst sichtbar gehaltene Unsicherheit. Iteration A soll den Systemzweck, den Kontext, die Systemgrenze und die wichtigsten Randbedingungen klaeren. Sie soll jedoch noch keine detaillierte Loesungsarchitektur, keine vollstaendige Sicherheitsauslegung und keine endgueltige Komponentenwahl erzwingen.

Offene Punkte werden deshalb in Anhang~\ref{chap:open-points} als eigenes Register gefuehrt. Dort ist dokumentiert:
\begin{itemize}
  \item welcher Punkt offen ist,
  \item warum er in Iteration A noch nicht entschieden wird,
  \item welche spaetere Iteration ihn voraussichtlich aufloest,
  \item welche Entwicklungsbereiche betroffen sind,
  \item welche Prioritaet der Punkt fuer die weitere Arbeit hat.
\end{itemize}

Diese Trennung schuetzt die Baseline A.1 vor zwei typischen Fehlern: Erstens werden offene Fragen nicht versehentlich als Anforderungen behandelt. Zweitens werden technische Loesungen nicht zu frueh festgelegt, bevor Use Cases, Requirements, Safety-Betrachtung und Verifikationskriterien ausreichend beschrieben sind.

\section{Uebergang zu Iteration B}

Iteration B beginnt nicht mit einer unstrukturierten Requirements-Liste. Stattdessen werden die Betriebsfaelle OP-01 bis OP-16 auf etwa acht bis zehn saubere Use Cases verdichtet. Aus diesen Use Cases werden funktionale Anforderungen, Leistungsanforderungen, Schnittstellenanforderungen, Safety-Anforderungen, Betriebsanforderungen und Verifikationsideen abgeleitet.

Die erwartete Arbeitsfolge fuer Iteration B ist:
\begin{enumerate}
  \item Betriebsfaelle priorisieren und gruppieren,
  \item Use Cases mit Vorbedingungen, Normalablauf, Alternativen und Ergebnissen beschreiben,
  \item Stakeholder Needs auf Requirements abbilden,
  \item Akzeptanzideen und Verifikationsmethoden ergaenzen,
  \item MVP-relevante Requirements markieren,
  \item offene Quantifizierungen fuer Werkstuecktraeger, Energie, Geschwindigkeitstoleranz, Sicherheitsabstand und Kartenlogik bearbeiten,
  \item Traceability-Matrix aktualisieren.
\end{enumerate}

\chapter{Iteration B: Use Cases und System Requirements}
\label{chap:iteration-b}

\section{Ziel der Iteration}

Iteration B ueberfuehrt die Systemdefinition aus Baseline A.1 in eine erste belastbare Use-Case- und Requirements-Baseline. Die Betriebsfaelle aus Iteration A werden verdichtet, beschrieben und als Quelle fuer testbare Anforderungen genutzt. Ziel ist noch nicht die Auswahl konkreter Hardware, Sensorik oder Softwarepakete, sondern die klaerere Beschreibung dessen, was das System leisten und nachweisen muss.

Die Leitfrage dieser Iteration lautet:
\begin{quote}
Welche beobachtbaren Systemleistungen muss der Mobile Material Robot im Zielsystem und im MVP erbringen, und wie kann spaeter nachgewiesen werden, dass diese Leistungen erfuellt sind?
\end{quote}

\section{Eingangsbasis aus Iteration A}

Iteration B basiert auf der \emph{System Definition Baseline A.1}. Besonders relevant sind:
\begin{itemize}
  \item Systemzweck: autonomer Transport von Werkstuecktraegern in einem produktionsaehnlichen Schulbereich,
  \item Systemgrenze: Roboter als System of Interest, Web-App, Ladestation, Werkstuecktraeger, Transfersystem und Umgebung als externe Schnittstellen,
  \item Betriebsumgebung: mehrere Raeume, Flur, offene Tueren, semi-dynamische Umgebung,
  \item Randbedingungen: ca. $0{,}5\,\mathrm{m/s}$, ca. 2 h Betriebsdauer, Web-App als erste Auftragsschnittstelle,
  \item Open-Issue-Register: offene Quantifizierungen und Architekturentscheidungen bleiben sichtbar.
\end{itemize}

\section{Verdichtung der Betriebsfaelle zu Use Cases}

Die 16 Betriebsfaelle aus Iteration A werden fuer Iteration B zu zehn Use Cases verdichtet. Dadurch bleiben die Use Cases ueberschaubar, decken aber weiterhin den End-to-End-Ablauf des Demonstrators ab.

\begin{longtable}{L{0.16\textwidth} L{0.28\textwidth} L{0.30\textwidth} L{0.14\textwidth}}
\toprule
Use Case & Name & Abgedeckte Betriebsfaelle & MVP \\
\midrule
UC-B-001 & Transportauftrag erteilen & OP-01 & Ja \\
UC-B-002 & Werkstuecktraeger bereitstellen & OP-02, OP-03 & Teilweise \\
UC-B-003 & Autonome Transportmission ausfuehren & OP-04, OP-05, OP-06, OP-07, OP-12 & Ja \\
UC-B-004 & Temporaeres Hindernis behandeln & OP-08, OP-09 & Ja \\
UC-B-005 & Blockierten Fahrweg behandeln & OP-10 & Ja \\
UC-B-006 & Person im Fahrbereich behandeln & OP-11, OP-16 & Ja, reduziert \\
UC-B-007 & Betrieblichen Stopp ausloesen & OP-13 & Ja \\
UC-B-008 & Rueckkehr zur Ladestation ausloesen & OP-14 & Teilweise \\
UC-B-009 & Technischen Fehler behandeln & OP-15, OP-16 & Teilweise \\
UC-B-010 & Karte und Umgebungsinformationen verwalten & OP-05, OP-08, OP-10 & Teilweise \\
\bottomrule
\end{longtable}

\section{Ausgearbeitete Kern-Use-Cases}

\subsection{UC-B-001: Transportauftrag erteilen}

\begin{longtable}{L{0.24\textwidth} L{0.64\textwidth}}
\toprule
Feld & Beschreibung \\
\midrule
Primaerer Akteur & Bediener ueber Web-App. \\
Ziel & Ein Transportauftrag mit Start- und Zielstandort wird eindeutig an den Roboter uebergeben. \\
Vorbedingungen & Web-App ist erreichbar; Roboter ist betriebsbereit oder meldet einen Zustand, in dem keine Annahme moeglich ist; Start und Ziel sind im System bekannt. \\
Normalablauf & Bediener waehlt Start und Ziel; Web-App sendet Auftrag; Roboter validiert Auftrag; Roboter meldet Annahme oder Ablehnung; Missionsstatus wird angezeigt. \\
Alternativen & Ziel unbekannt; Roboter nicht bereit; Ladestand unzureichend; Fahrbereich gesperrt; Auftrag wird abgelehnt und begruendet gemeldet. \\
Ergebnis & Auftrag ist angenommen und besitzt einen Status oder wurde nachvollziehbar abgelehnt. \\
Abgeleitete Requirements & REQ-B-MIS-001, REQ-B-UI-001, REQ-B-STA-001. \\
\bottomrule
\end{longtable}

\subsection{UC-B-003: Autonome Transportmission ausfuehren}

\begin{longtable}{L{0.24\textwidth} L{0.64\textwidth}}
\toprule
Feld & Beschreibung \\
\midrule
Primaerer Akteur & Bediener als Auftraggeber; Roboter als autonom ausfuehrendes System. \\
Ziel & Der Roboter faehrt autonom vom Startbereich zum Zielbereich und erreicht dabei mehrere Raeume ueber Flur und offene Tueren. \\
Vorbedingungen & Auftrag ist angenommen; Grundkarte ist geladen; Roboter ist lokalisiert; Ziel ist bekannt; Fahrbereich ist grundsaetzlich befahrbar. \\
Normalablauf & Roboter plant eine Route; faehrt mit MVP-Nominalgeschwindigkeit; durchfaehrt offene Tueren und Flurabschnitte; aktualisiert Missionsstatus; meldet Ziel erreicht. \\
Alternativen & Hindernis wird erkannt; Person befindet sich im Fahrbereich; Lokalisierung unsicher; Ziel nicht erreichbar; Roboter stoppt, plant neu oder meldet Blockade. \\
Ergebnis & Ziel wird erreicht oder die Nichtausfuehrung wird mit Grund gemeldet. \\
Abgeleitete Requirements & REQ-B-NAV-001, REQ-B-NAV-002, REQ-B-NAV-003, REQ-B-PERF-001, REQ-B-LOC-001, REQ-B-STA-001. \\
\bottomrule
\end{longtable}

\subsection{UC-B-004: Temporaeres Hindernis behandeln}

\begin{longtable}{L{0.24\textwidth} L{0.64\textwidth}}
\toprule
Feld & Beschreibung \\
\midrule
Primaerer Akteur & Umgebung als Ereignisquelle; indirekt Bediener und Personen im Fahrbereich. \\
Ziel & Der Roboter verhindert Kollisionen und setzt die Mission fort, falls dies sicher moeglich ist. \\
Vorbedingungen & Roboter befindet sich in Fahrt oder Planung; Sensorik liefert aktuelle Umgebungsinformationen. \\
Normalablauf & Hindernis wird erkannt; lokale Umgebung wird aktualisiert; Roboter stoppt, weicht aus oder plant lokal neu; Mission wird fortgesetzt. \\
Alternativen & Hindernis ist zu nah; Hindernis blockiert alle sinnvollen Wege; Wahrnehmung ist unsicher; Roboter geht in sicheren Zustand und meldet Blockade. \\
Ergebnis & Keine Kollision; Mission fortgesetzt oder kontrolliert blockiert. \\
Abgeleitete Requirements & REQ-B-PER-001, REQ-B-NAV-004, REQ-B-SAFE-001, REQ-B-MAP-002. \\
\bottomrule
\end{longtable}

\subsection{UC-B-006: Person im Fahrbereich behandeln}

\begin{longtable}{L{0.24\textwidth} L{0.64\textwidth}}
\toprule
Feld & Beschreibung \\
\midrule
Primaerer Akteur & Person im Fahrbereich. \\
Ziel & Personen im Fahrbereich werden bei der Bewegungsentscheidung sicherheitsrelevant beruecksichtigt. \\
Vorbedingungen & Roboter faehrt oder plant eine Bewegung; Personen koennen den Fahrbereich betreten. \\
Normalablauf & Roboter erkennt eine sicherheitsrelevante Situation; Bewegung wird reduziert, gestoppt oder umgeplant; Status wird gemeldet. \\
Alternativen & Erkennung unsicher; Person bleibt im Fahrbereich; Situation ist nicht sicher aufloesbar; Roboter bleibt im sicheren Zustand. \\
Ergebnis & Keine gefaehrdende Bewegung gegenueber Personen; Zustand ist nachvollziehbar. \\
Abgeleitete Requirements & REQ-B-SAFE-002, REQ-B-SAFE-003, REQ-B-STA-001. \\
\bottomrule
\end{longtable}

\section{Requirement-Baseline B.1}

Die folgende Tabelle bildet eine erste Requirements-Baseline B.1. Sie ersetzt die Anforderungen aus Iteration A nicht, sondern praezisiert sie und macht ihre Quelle, Akzeptanzidee und Verifikation sichtbarer.

\begin{longtable}{L{0.17\textwidth} L{0.14\textwidth} L{0.32\textwidth} L{0.23\textwidth}}
\toprule
ID & Typ & Requirement & Akzeptanz- und Verifikationsidee \\
\midrule
\modeltable{requirements_b1}
\bottomrule
\end{longtable}

\section{MVP-Markierung der Iteration B}

\begin{longtable}{L{0.20\textwidth} L{0.16\textwidth} L{0.50\textwidth}}
\toprule
Requirement & MVP-Status & Begruendung \\
\midrule
\modeltable{mvp_b1}
\bottomrule
\end{longtable}

\section{Open-Issue-Auswirkung}

Iteration B loest einige Open Issues teilweise auf, fuehrt andere aber bewusst weiter.

\begin{longtable}{L{0.16\textwidth} L{0.40\textwidth} L{0.28\textwidth}}
\toprule
Open Issue & Stand nach Iteration B & Naechste Aktion \\
\midrule
OI-MECH-001 & Weiter offen. Requirement REQ-B-TRN-001 macht die Payload-Faehigkeit sichtbar, quantifiziert sie aber nicht. & Daten des Transfersystems beschaffen. \\
OI-MECH-002 & Weiter offen. Use Case UC-B-002 beschreibt die Bereitstellung, nicht die mechanische Konstruktion. & Mechanische Schnittstellenanforderungen ableiten. \\
OI-NAV-001 & Teilweise bearbeitet. REQ-B-MAP-001 fordert eine freigegebene Grundkarte. & Kartenprozess in Architekturphase definieren. \\
OI-NAV-002 & Teilweise bearbeitet. REQ-B-MAP-002 fordert die Trennung temporaer/dauerhaft. & Regeln fuer Kartenupdate-Freigabe modellieren. \\
OI-SAFE-001 & Weiter offen. REQ-B-SAFE-003 trennt Begriffe, ersetzt aber kein Safety-Konzept. & Safety-Konzept in Iteration D/F erstellen. \\
OI-SAFE-002 & Weiter offen. REQ-B-PERF-001 haelt $0{,}5\,\mathrm{m/s}$ fest, Toleranzen und Bremsweg fehlen. & Mess- und Safety-Annahmen festlegen. \\
OI-ENE-001 & Teilweise bearbeitet. REQ-B-PWR-001 fordert ein 2-h-Unterrichtsprofil. & Profil quantitativ definieren. \\
OI-ENE-002 & Teilweise bearbeitet. REQ-B-ENE-001 deckt Rueckkehr als Fahrziel ab. & Docking/Laden im MVP-Schnitt entscheiden. \\
OI-UI-001 & Teilweise bearbeitet. UC-B-001 und REQ-B-UI-001 legen Auftragserteilung fest. & UI-Rollen und Diagnoseumfang spezifizieren. \\
OI-TEST-001 & Vorbereitet. Requirements besitzen Verifikationsideen. & Testumgebung in Iteration F ausarbeiten. \\
\bottomrule
\end{longtable}

\section{Definition of Done fuer Iteration B}

Iteration B gilt fuer diese Arbeitsfassung als abgeschlossen, wenn:
\begin{itemize}
  \item die Betriebsfaelle aus Iteration A auf eine ueberschaubare Use-Case-Struktur verdichtet sind,
  \item die wichtigsten Use Cases mit Vorbedingungen, Normalablaeufen, Alternativen und Ergebnissen beschrieben sind,
  \item jedes zentrale MVP-Verhalten mindestens eine Candidate Requirement besitzt,
  \item jede Candidate Requirement eine Akzeptanz- oder Verifikationsidee hat,
  \item Open Issues nicht verschwinden, sondern mit Requirements, Entscheidungen oder spaeteren Aktionen verknuepft sind.
\end{itemize}

Die naechste Iteration ist Iteration C. Sie sollte den formalen MVP-Schnitt festlegen: welche Requirements in MVP 1 verpflichtend sind, welche nur als Schnittstelle sichtbar bleiben und welche bewusst auf spaetere Releases verschoben werden.

\chapter{Glossar}
\label{chap:glossar}

Dieses Glossar erklaert Abkuerzungen und Fachbegriffe, die in der Dokumentation verwendet werden. Es ist bewusst als normales Kapitel umgesetzt, damit das Projekt ohne zusaetzlichen Glossar-Build-Schritt in Overleaf kompilierbar bleibt.

\section{Abkuerzungen}

\begin{longtable}{L{0.20\textwidth} L{0.30\textwidth} L{0.36\textwidth}}
\toprule
Abkuerzung & Langform & Bedeutung im Projekt \\
\midrule
ADR & Architecture Decision Record & Dokumentierte Architektur- oder Prozessentscheidung mit Kontext, Alternativen, Begruendung und Auswirkungen. \\
AI / KI & Artificial Intelligence / Kuenstliche Intelligenz & Hilfsmittel zur Erstellung, Pflege und Pruefung von Entwicklungsartefakten; keine alleinige Freigabeinstanz. \\
HW & Hardware & Physische Bestandteile wie Mechanik, Sensorik, Aktoren, Energieversorgung und Rechner. \\
IF & Interface & Eindeutig benannte Schnittstelle zwischen Systemelementen oder externen Akteuren. \\
MBSE & Model-Based Systems Engineering & Modellbasierter Entwicklungsansatz, bei dem Anforderungen, Funktionen, Architektur, Schnittstellen und Tests verbunden werden. \\
MMR & Mobile Material Robot & Kurzname des betrachteten mobilen Materialtransportroboters. \\
MVP & Minimum Viable Product & Kleinster durchgaengiger Systemausschnitt, der die zentrale technische Hypothese nachweist. \\
OP & Operating Case / Betriebsfall & Reale Betriebssituation als Vorstufe zu einem formalisierten Use Case. \\
REQ & Requirement & Nachzuweisende Anforderung an das System. \\
ROS2 & Robot Operating System 2 & Wahrscheinliche spaetere Robotik-Middleware; in Iteration A noch keine feste technische Anforderung. \\
SN & Stakeholder Need & Beduerfnis oder Erwartung eines Stakeholders, noch nicht zwingend testbar formuliert. \\
SoI & System of Interest & Das im Projekt betrachtete und zu entwickelnde System. \\
SW & Software & Programme, Konfigurationen und Laufzeitlogik des Roboters. \\
TC & Test Case & Testfall zur Verifikation einer oder mehrerer Anforderungen. \\
UC & Use Case & Zielorientierte Beschreibung einer Interaktion zwischen Akteur und System. \\
\bottomrule
\end{longtable}

\section{Fachbegriffe}

\begin{longtable}{L{0.24\textwidth} L{0.62\textwidth}}
\toprule
Begriff & Erklaerung \\
\midrule
Anforderung & Aussage darueber, was das System leisten, einhalten oder nachweisen muss. Anforderungen sollen notwendig, eindeutig, pruefbar und moeglichst loesungsneutral sein. \\
Annahme & Vorlaeufige Aussage, die fuer die Weiterarbeit verwendet wird, aber spaeter bestaetigt, geaendert oder durch eine Entscheidung ersetzt werden muss. \\
Baseline & Festgehaltener, referenzierbarer Arbeitsstand eines Entwicklungsartefakts. Baseline A.1 ist der freigegebene Stand der Systemdefinition nach Iteration A. \\
Betrieblicher Stopp & Bedienbefehl zum kontrollierten Stoppen oder Pausieren des Betriebs. Er ist nicht gleichbedeutend mit einem sicherheitsgerichteten Not-Halt. \\
Demonstrator & System, das nicht nur eine technische Funktion ausfuehrt, sondern auch Lern-, Demonstrations- und Erklaerzwecke erfuellt. \\
EmergencyStop / Not-Halt & Sicherheitsgerichtete Funktion zum schnellen Erreichen eines sicheren Zustands. Bedarf, Ausfuehrung und Normbezug werden spaeter in der Sicherheitsbetrachtung festgelegt. \\
Grundkarte & Vorab bekannte Karte des Betriebsbereichs, zum Beispiel Raeume, Flure, feste Durchgaenge und bekannte Transportstandorte. \\
Hindernis & Objekt oder Person im Fahrbereich, das die geplante Bewegung beeinflusst oder verhindert. \\
Langzeitkarte & Dauerhaft gespeicherte Karteninformation. Sie darf nicht unkontrolliert durch temporaere Hindernisse veraendert werden. \\
Lokalisierung & Bestimmung der Roboterposition und -orientierung innerhalb einer Karte oder eines Koordinatensystems. \\
Mission & Logische Ausfuehrung eines Transportauftrags vom Annehmen des Auftrags bis zum Abschluss, Abbruch oder Fehlerzustand. \\
Nicht-Ziel & Bewusst ausgeschlossener Umfang, der hilft, den Scope einer Iteration zu begrenzen. \\
Payload & Nutzlast des Roboters; in diesem Projekt der Werkstuecktraeger einschliesslich eventuell transportierter Werkstuecke. \\
Perception & Wahrnehmungsfunktion des Roboters, die Sensordaten interpretiert und fuer Navigation, Karte oder Safety nutzbar macht. \\
SafetyStop & Sicherheitsbezogene Reaktion auf eine erkannte unsichere Situation. Die genaue Ausfuehrung ist Teil der spaeteren Safety-Architektur. \\
Semi-dynamische Umgebung & Umgebung mit teilweise stabiler Struktur, aber veraenderlichen Hindernissen, Personen oder Gegenstaenden. \\
Schnittstelle & Vereinbarter Beruehrungspunkt zwischen Systemelementen oder externen Systemen; kann mechanisch, elektrisch, informationstechnisch, energetisch oder organisatorisch sein. \\
Systemgrenze & Gedachte Grenze, die festlegt, was zum System gehoert und was als externes System, Akteur oder Umgebung behandelt wird. \\
Traceability & Nachvollziehbare Verknuepfung von Needs, Requirements, Use Cases, Funktionen, Architektur, Schnittstellen, Implementierung und Tests. \\
Transfersystem & Externes System, das Werkstuecktraeger in einer produktionsaehnlichen Umgebung bewegt oder bereitstellt. \\
Use Case & Strukturierte Beschreibung eines Ziels, eines Akteurs, der Vorbedingungen, des Normalablaufs, der Alternativen und des Ergebnisses. \\
Validierung & Pruefung, ob das entwickelte System das tatsaechliche Problem und die Beduerfnisse der Stakeholder trifft. \\
Verifikation & Pruefung, ob das System die spezifizierten Anforderungen erfuellt. \\
Werkstuecktraeger & Standardisierter Traeger eines Transfersystems, der Werkstuecke aufnehmen kann und im Projekt als Transportgut dient. \\
\bottomrule
\end{longtable}


\appendix
\chapter{Entscheidungsprotokoll}

\section{Entscheidungen aus Baseline A.1}

\begin{longtable}{L{0.14\textwidth} L{0.42\textwidth} L{0.14\textwidth} L{0.20\textwidth}}
\toprule
ID & Entscheidung & Status & Begruendung \\
\midrule
\modeltable{decisions_baseline}
\bottomrule
\end{longtable}

\section{Architektur- und Prozessentscheidungen}

\begin{longtable}{L{0.14\textwidth} L{0.22\textwidth} L{0.22\textwidth} L{0.26\textwidth}}
\toprule
ID & Entscheidung & Status & Begruendung \\
\midrule
\modeltable{decisions_architecture}
\bottomrule
\end{longtable}

\section{ADR-Template}

Jede neue Architekturentscheidung wird in dieser Struktur dokumentiert:
\begin{itemize}
  \item Kontext,
  \item Entscheidung,
  \item betrachtete Alternativen,
  \item Begruendung,
  \item Auswirkungen auf Zielsystem,
  \item Auswirkungen auf MVP,
  \item Risiken,
  \item betroffene Requirements,
  \item betroffene Tests,
  \item Datum und Status.
\end{itemize}

\chapter{Open-Issue-Register}
\label{chap:open-points}

\section{Zweck des Registers}

Dieses Register listet alle Punkte, die in der \emph{System Definition Baseline A.1} bewusst offen bleiben. Die offenen Punkte sind nicht aus Versehen ungeklaert. Sie werden zurueckgestellt, weil Iteration A nur den Systemzweck, den Kontext, die Systemgrenze, zentrale Randbedingungen und die fachliche Ausgangsbasis festlegt.

Die weitere Bearbeitung erfolgt in spaeteren Iterationen:
\begin{itemize}
  \item Iteration B: Use Cases und testbare System Requirements,
  \item Iteration C: formaler MVP-Schnitt und Erfolgskriterien,
  \item Iteration D: logische und physische Systemarchitektur,
  \item Iteration E: HW/SW-Allokation, ROS2, Elektronik und Mechanik,
  \item Iteration F: Verifikation, Testumgebung und Traceability-Abschluss.
\end{itemize}

\section{Open Issues nach Baseline A.1}

\begin{longtable}{L{0.13\textwidth} L{0.30\textwidth} L{0.22\textwidth} L{0.17\textwidth} L{0.08\textwidth}}
\toprule
ID & Offener Punkt & Warum zurueckgestellt & Geplante Aufloesung & Prio. \\
\midrule
OI-MECH-001 & Konkretes Transfersystem, Abmessungen, Eigenmasse, maximale Werkstueckmasse, Gesamtmasse und Schwerpunktbereich des Werkstuecktraegers festlegen. & Diese Werte haengen vom realen Transfersystem ab und sollen nicht geraten werden. & Iteration B/C: Requirements und MVP-Schnitt; Iteration E: Mechanik. & Hoch \\
OI-MECH-002 & Mechanische Transportaufnahme definieren: Auflage, Fixierung, Toleranzen, Verrutschen, manuelles Handling. & Die Aufnahme muss aus Payload-Daten, Safety-Betrachtung und Use Cases abgeleitet werden. & Iteration D/E: Architektur und Mechanik. & Hoch \\
OI-NAV-001 & Kartenerstellung, Kartenfreigabe und Versionierung fuer die Grundkarte definieren. & Iteration A trennt nur die Informationsarten; der Kartenprozess braucht Requirements und Betriebskonzept. & Iteration B/D: Requirements und Navigationsarchitektur. & Hoch \\
OI-NAV-002 & Regeln festlegen, wann eine erkannte Aenderung temporaer bleibt oder als dauerhaftes Kartenupdate markiert wird. & Eine vorschnelle Festlegung koennte falsche Kartenupdates oder zu starre Navigation erzeugen. & Iteration B/D/F: Requirements, Architektur, Tests. & Hoch \\
OI-SAFE-001 & Sicherheitskonzept fuer Personen im Fahrbereich, SafetyStop und EmergencyStop klaeren. & SafetyStop, EmergencyStop und betrieblicher Stopp muessen normativ und technisch getrennt betrachtet werden. & Iteration B/D/F: Safety-Requirements, Architektur, Verifikation. & Hoch \\
OI-SAFE-002 & Mindestabstaende, Bremsweg, Hindernisdetektionsdistanz und Geschwindigkeitstoleranz fuer $0{,}5\,\mathrm{m/s}$ festlegen. & Messwerte haengen von Masse, Bremsverhalten, Sensorik, Regelung und Safety-Konzept ab. & Iteration B/F: quantifizierte Requirements und Tests. & Hoch \\
OI-ENE-001 & Nutzungsprofil fuer den 2-h-Unterrichtsbetrieb definieren. & Die 2 h sind als Ziel festgelegt; das pruefbare Lastprofil muss noch beschrieben werden. & Iteration B/F: Energie-Requirement und Testprofil. & Mittel \\
OI-ENE-002 & Entscheiden, ob Rueckkehr zur Ladestation nur Anfahren, automatisches Docking oder automatisches Laden umfasst. & Die Rueckkehrfunktion ist Zielsystembestandteil, aber der MVP-Umfang muss noch geschnitten werden. & Iteration C/D: MVP-Scope und Architektur. & Mittel \\
OI-UI-001 & Informationsumfang und Rollen der Web-App bestimmen: Auftrag, Status, Stopp, Rueckkehr, Fehler, Diagnose. & Iteration A legt die Web-App als Schnittstelle fest, aber nicht alle UI-Funktionen. & Iteration B/C: Use Cases und MVP-Scope. & Mittel \\
OI-IT-001 & Netzwerk, Zugriffsschutz, Datenhaltung und Update-Konzept fuer den Demonstrator klaeren. & IT-Details sollen aus Betrieb, Datenschutz, Unterrichtsnutzung und Web-App-Architektur folgen. & Iteration D/E: Kommunikations- und IT-Architektur. & Mittel \\
OI-TEST-001 & Testumgebung fuer Raeume, Flur, offene Tuer, Hindernisse und Personennaehe definieren. & Testaufbau und Kriterien benoetigen quantifizierte Requirements und Safety-Annahmen. & Iteration F, vorbereitet in B/C. & Hoch \\
\bottomrule
\end{longtable}

\section{Bewusst nicht offene Punkte}

Folgende Punkte gelten fuer Baseline A.1 als entschieden und werden nicht mehr als offene Abschlussfragen der Iteration A gefuehrt:
\begin{itemize}
  \item Einsatzort: produktionsaehnlicher Schulbereich,
  \item Fahrbereich: mehrere Raeume und Flur,
  \item offene Tueren muessen durchfahren werden,
  \item Transportgut: Werkstuecktraeger eines Transfersystems,
  \item Beladung: zunaechst extern durch Mensch oder Transfersystem,
  \item Geschwindigkeit MVP: ca. $0{,}5\,\mathrm{m/s}$,
  \item Betriebsdauer erste Version: ca. 2 h,
  \item Auftragserteilung erste Version: Web-App.
\end{itemize}

\section{Pflegeregel}

Jedes Open Issue bleibt im Register, bis es durch eine Anforderung, eine Entscheidung, einen Testfall oder eine bewusst dokumentierte Nicht-Ziel-Entscheidung geschlossen wird. Beim Schliessen wird die betroffene Requirement-, Decision- oder Test-ID referenziert. Dadurch bleibt nachvollziehbar, warum ein Punkt geloest wurde und welche Artefakte davon betroffen sind.

\chapter{Initiale Traceability-Matrix}

\begin{longtable}{L{0.15\textwidth} L{0.18\textwidth} L{0.20\textwidth} L{0.20\textwidth} L{0.14\textwidth}}
\toprule
Requirement & Need/Funktion & Systemelement & Schnittstelle/ROS2 & Test \\
\midrule
\modeltable{traceability_a1_initial}
\bottomrule
\end{longtable}

\section{Ergaenzte Traceability aus Iteration A}

\begin{longtable}{L{0.15\textwidth} L{0.18\textwidth} L{0.20\textwidth} L{0.20\textwidth} L{0.14\textwidth}}
\toprule
Requirement & Need/Funktion & Systemelement & Schnittstelle/ROS2 & Test \\
\midrule
\modeltable{traceability_a1_iteration_a}
\bottomrule
\end{longtable}


\printbibliography

\end{document}
